\documentclass[11pt]{amsart}
\usepackage[margin=1in]{geometry}
\usepackage{amsmath,amssymb,amsthm,mathtools}
\usepackage{amscd}
\usepackage[hidelinks]{hyperref}

% Zenodo / DOI companion layout (shared by closure, appendix, long-form drafts).
% The file ``companion-code.zip'' is published alongside this record under the same DOI.
% After unpacking, the Lean/scripts/certificate mirror for this manuscript is
% ``release/paper_refs_bundle_2026-05-06/'' (see ``release/paper_refs_bundle_2026-05-06/MANIFEST.json'').
\newcommand{\CompanionCodeZip}{\path{companion-code.zip}}
\newcommand{\PaperReproBundleDir}{\path{release/paper_refs_bundle_2026-05-06}}
\newcommand{\PaperReproBundleManifest}{\path{release/paper_refs_bundle_2026-05-06/MANIFEST.json}}
\newcommand{\ReleaseDistArchive}{\CompanionCodeZip}

% Back-compat names used throughout this long-form draft:
\newcommand{\PaperArtifactBundleZip}{\ReleaseDistArchive}
\newcommand{\PaperArtifactBundleDir}{\PaperReproBundleDir}

\title[Local Rapidity Fiber to Spin(8)]{From Local Rapidity Fiber to Spin(8):\\
The Primitive Rapidity Observable and 28-Dimensional Closure}
\author{Steven Ettinger Jr}
\address{Independent Researcher}
\email{steven@disregardfiat.tech}
\date{Preprint v1 --- \today}

\newtheorem{theorem}{Theorem}[section]
\newtheorem{lemma}[theorem]{Lemma}
\newtheorem{proposition}[theorem]{Proposition}
\newtheorem{definition}[theorem]{Definition}
\newtheorem{remark}[theorem]{Remark}
\newtheorem{example}[theorem]{Example}

% Reproducibility: see ``include/release_archive_macros.tex''; inner mirror is \PaperReproBundleDir/.

\begin{document}
\begin{abstract}
We present a single derivation chain on a discrete index space. Under a uniform polynomial growth hypothesis, a positive, divergent cumulative curvature channel \(K\) is obtained (for any fixed \(\alpha>0\) in the density \(\rho\); the manuscript highlights \(\alpha=3/5\) as the HQIV/Brodie normalization). Normalizing \(K\) yields a \emph{normalized cumulative phase readout} \(\mathcal{R}\) built from explicit phase-map axioms; we use the name ``rapidity'' only as an interpretive bridge to hyperbolic kinematics and geometric-phase language. The geometric phase direction extracted from this readout supplies a concrete antisymmetric phase-lift generator \(\Delta\) on a \emph{chosen} \(2\)-plane in \(\mathbb{R}^8\). Once the standard octonionic derivation matrices \(G_2\) and this HQVM/Fano-aligned \(\Delta\) are fixed as external mathematical inputs, adjoining \(\Delta\) to the \(14\) generators of \(G_2\) closes, in the Lie sense, to the full \(28\)-dimensional algebra \(\mathfrak{so}(8)\) (machine-checked). Causal set theory enters only as motivation: if a discrete model's shell increment matches the same quadratic law used here, the \emph{same} analytic channel and the \emph{same} Lie data apply; causal partial order alone does \emph{not} select the octonion multiplication table, the color axis \(e_7\), or the \((e_1,e_7)\) phase plane. The same construction generalizes to any \(\mathrm{SO}(n)\) via the Fourier-twiddle route on \(S^{n-1}\). Historically this closure was first obtained through the full rapidity fiber; the concise \(G_2+\Delta\) route is its algebraic distillation. All nontrivial closure, transport, and holonomy steps cited from the Lean library are machine-checked there under the hypotheses stated in the text. For the full symbolic certificate pipeline as a minimal standalone PDF, see \path{papers/so8_closure_full_appendix.pdf} (source \path{papers/so8_closure_full_appendix.tex}); those sources appear under \PaperReproBundleDir/ within the release distribution \ReleaseDistArchive{} archived with the submission rewrite \path{papers/closure.tex} (see \PaperReproBundleManifest). \textbf{Note:} this file is the long-form workspace manuscript; \path{papers/closure.tex} is the standalone rewrite used for publication.
\end{abstract}
\maketitle

% Shared reader contract: patch QFT + discrete numerics.
% From papers/*.tex:     % Shared reader contract: patch QFT + discrete numerics.
% From papers/*.tex:     % Shared reader contract: patch QFT + discrete numerics.
% From papers/*.tex:     \input{include/patch_theory_messaging.tex}
% From papers/paper/*.tex: \input{../include/patch_theory_messaging.tex}

\paragraph{Patch QFT and discrete numerics (reader contract).}
HQIV (Horizon-Quantized Informational Vacuum) is formulated as a \emph{patch theory}: quantum fields, actions, and physical readouts are attached to discrete patches---shell indices $m\in\mathbb{N}$, finite spacetime index types (e.g.\ $\mathrm{Fin}\,4$), and horizon-limited charts---rather than to a hypothetical fundamental smooth spacetime obtained as a mesh-refinement or ultraviolet limit.
Accordingly, when this text refers to ``QFT,'' it means \emph{patch QFT}: patching rules, finite measurement ledgers, and variational identities on the discrete spine; it is not the claim that the theory is \emph{defined} by recovering ordinary continuum quantum field theory through such a limit.
The numbers that admit the sharpest statements---mass and coupling ladders, curvature ratios, shell thermodynamics, and rational witnesses in \texttt{HQIV\_LEAN}---are objects of \emph{discrete mathematics} on $\mathbb{N}$ (combinatorial growth, cumulative simplex ledgers) together with the ordered field $\mathbb{R}$ as used in the proof assistant for analysis, bounds, and phenomenological display.
Smooth-manifold or continuum field notation, when it appears, is a \emph{translation layer} for comparison with standard physics literature; it does not supersede the discrete definitions.

\paragraph{A continuum is not required, and in places would destroy these results.}
The discrete patch layer is \emph{not} a numerical approximation to a more fundamental continuum theory awaiting future construction; it is the ontology.  Where the literature treats ``the continuum limit'' as the goal, HQIV treats it as a \emph{calculation approximation} for comparison with standard formulas.  The continuum form is an optional convenience for comparison; the proved discrete statement is what carries the physics.  In several load-bearing places---the curvature imprint $\delta_E(m)$ at shell $m$, the harmonic ladder masses $m_\tau\to m_\mu\to m_e$, the lock-in vev $v=\sqrt{\eta_{\rm paper}\,\Omega_k(m_{\rm lockin};m_{\rm lockin})}$, the proton anchor at $m=\mathrm{referenceM}=4$, the baryon-asymmetry $\eta\propto 6^7\sqrt{3}/(m+1)$, and the rational coefficients $\alpha=3/5$, $\gamma=2/5$, $\alpha_{\rm GUT}=1/42$---taking a literal $m\to\infty$ or sub-Planck continuum limit would \emph{erase} the discrete index that is doing the predictive work.  Continuum forms of these objects are therefore not preferred refinements; they are degenerate, information-losing readouts of the same patch data.

\paragraph{An exception: $\mathfrak{so}(8)$ as a de-facto continuum algebra from a discrete construction.}
Principled exception to the discrete-only stance: the closed Lie algebra $\mathfrak{so}(8) \supseteq G_2\cup\{\Delta\}$ generated from the Fano octonion left-multiplication table and the phase-lift generator $\Delta\in(\mathrm{e}_1,\mathrm{e}_7)$.  Although $\mathfrak{so}(8)$ is a continuous (28-dimensional real) Lie algebra, it is built here entirely from \emph{discrete} ingredients (the seven imaginary octonion units, a finite Fano table, and one $\pi/2$ rotation), and its closure to 28 dimensions is a finite linear-algebra certificate (see the \texttt{Hqiv.SO8Closure} / \texttt{Hqiv.SO8ClosureInterface} stack).  The Lie-group structure is therefore a \emph{de-facto continuum algebra obtained from a discrete construction}: continuous in parameter space, but with no continuum spacetime, no continuum measure, and no UV limit attached.  When this manuscript or its companions write $\exp(t\,\Delta)\in\mathrm{SO}(8)$ or treat $\mathfrak{so}(8)$ rotations as smooth, those are statements internal to a finite-dimensional Lie group whose generators are certified by discrete data; they are \emph{not} a continuum-spacetime commitment and do not require a continuum-QFT scaffolding.

% From papers/paper/*.tex: % Shared reader contract: patch QFT + discrete numerics.
% From papers/*.tex:     \input{include/patch_theory_messaging.tex}
% From papers/paper/*.tex: \input{../include/patch_theory_messaging.tex}

\paragraph{Patch QFT and discrete numerics (reader contract).}
HQIV (Horizon-Quantized Informational Vacuum) is formulated as a \emph{patch theory}: quantum fields, actions, and physical readouts are attached to discrete patches---shell indices $m\in\mathbb{N}$, finite spacetime index types (e.g.\ $\mathrm{Fin}\,4$), and horizon-limited charts---rather than to a hypothetical fundamental smooth spacetime obtained as a mesh-refinement or ultraviolet limit.
Accordingly, when this text refers to ``QFT,'' it means \emph{patch QFT}: patching rules, finite measurement ledgers, and variational identities on the discrete spine; it is not the claim that the theory is \emph{defined} by recovering ordinary continuum quantum field theory through such a limit.
The numbers that admit the sharpest statements---mass and coupling ladders, curvature ratios, shell thermodynamics, and rational witnesses in \texttt{HQIV\_LEAN}---are objects of \emph{discrete mathematics} on $\mathbb{N}$ (combinatorial growth, cumulative simplex ledgers) together with the ordered field $\mathbb{R}$ as used in the proof assistant for analysis, bounds, and phenomenological display.
Smooth-manifold or continuum field notation, when it appears, is a \emph{translation layer} for comparison with standard physics literature; it does not supersede the discrete definitions.

\paragraph{A continuum is not required, and in places would destroy these results.}
The discrete patch layer is \emph{not} a numerical approximation to a more fundamental continuum theory awaiting future construction; it is the ontology.  Where the literature treats ``the continuum limit'' as the goal, HQIV treats it as a \emph{calculation approximation} for comparison with standard formulas.  The continuum form is an optional convenience for comparison; the proved discrete statement is what carries the physics.  In several load-bearing places---the curvature imprint $\delta_E(m)$ at shell $m$, the harmonic ladder masses $m_\tau\to m_\mu\to m_e$, the lock-in vev $v=\sqrt{\eta_{\rm paper}\,\Omega_k(m_{\rm lockin};m_{\rm lockin})}$, the proton anchor at $m=\mathrm{referenceM}=4$, the baryon-asymmetry $\eta\propto 6^7\sqrt{3}/(m+1)$, and the rational coefficients $\alpha=3/5$, $\gamma=2/5$, $\alpha_{\rm GUT}=1/42$---taking a literal $m\to\infty$ or sub-Planck continuum limit would \emph{erase} the discrete index that is doing the predictive work.  Continuum forms of these objects are therefore not preferred refinements; they are degenerate, information-losing readouts of the same patch data.

\paragraph{An exception: $\mathfrak{so}(8)$ as a de-facto continuum algebra from a discrete construction.}
Principled exception to the discrete-only stance: the closed Lie algebra $\mathfrak{so}(8) \supseteq G_2\cup\{\Delta\}$ generated from the Fano octonion left-multiplication table and the phase-lift generator $\Delta\in(\mathrm{e}_1,\mathrm{e}_7)$.  Although $\mathfrak{so}(8)$ is a continuous (28-dimensional real) Lie algebra, it is built here entirely from \emph{discrete} ingredients (the seven imaginary octonion units, a finite Fano table, and one $\pi/2$ rotation), and its closure to 28 dimensions is a finite linear-algebra certificate (see the \texttt{Hqiv.SO8Closure} / \texttt{Hqiv.SO8ClosureInterface} stack).  The Lie-group structure is therefore a \emph{de-facto continuum algebra obtained from a discrete construction}: continuous in parameter space, but with no continuum spacetime, no continuum measure, and no UV limit attached.  When this manuscript or its companions write $\exp(t\,\Delta)\in\mathrm{SO}(8)$ or treat $\mathfrak{so}(8)$ rotations as smooth, those are statements internal to a finite-dimensional Lie group whose generators are certified by discrete data; they are \emph{not} a continuum-spacetime commitment and do not require a continuum-QFT scaffolding.


\paragraph{Patch QFT and discrete numerics (reader contract).}
HQIV (Horizon-Quantized Informational Vacuum) is formulated as a \emph{patch theory}: quantum fields, actions, and physical readouts are attached to discrete patches---shell indices $m\in\mathbb{N}$, finite spacetime index types (e.g.\ $\mathrm{Fin}\,4$), and horizon-limited charts---rather than to a hypothetical fundamental smooth spacetime obtained as a mesh-refinement or ultraviolet limit.
Accordingly, when this text refers to ``QFT,'' it means \emph{patch QFT}: patching rules, finite measurement ledgers, and variational identities on the discrete spine; it is not the claim that the theory is \emph{defined} by recovering ordinary continuum quantum field theory through such a limit.
The numbers that admit the sharpest statements---mass and coupling ladders, curvature ratios, shell thermodynamics, and rational witnesses in \texttt{HQIV\_LEAN}---are objects of \emph{discrete mathematics} on $\mathbb{N}$ (combinatorial growth, cumulative simplex ledgers) together with the ordered field $\mathbb{R}$ as used in the proof assistant for analysis, bounds, and phenomenological display.
Smooth-manifold or continuum field notation, when it appears, is a \emph{translation layer} for comparison with standard physics literature; it does not supersede the discrete definitions.

\paragraph{A continuum is not required, and in places would destroy these results.}
The discrete patch layer is \emph{not} a numerical approximation to a more fundamental continuum theory awaiting future construction; it is the ontology.  Where the literature treats ``the continuum limit'' as the goal, HQIV treats it as a \emph{calculation approximation} for comparison with standard formulas.  The continuum form is an optional convenience for comparison; the proved discrete statement is what carries the physics.  In several load-bearing places---the curvature imprint $\delta_E(m)$ at shell $m$, the harmonic ladder masses $m_\tau\to m_\mu\to m_e$, the lock-in vev $v=\sqrt{\eta_{\rm paper}\,\Omega_k(m_{\rm lockin};m_{\rm lockin})}$, the proton anchor at $m=\mathrm{referenceM}=4$, the baryon-asymmetry $\eta\propto 6^7\sqrt{3}/(m+1)$, and the rational coefficients $\alpha=3/5$, $\gamma=2/5$, $\alpha_{\rm GUT}=1/42$---taking a literal $m\to\infty$ or sub-Planck continuum limit would \emph{erase} the discrete index that is doing the predictive work.  Continuum forms of these objects are therefore not preferred refinements; they are degenerate, information-losing readouts of the same patch data.

\paragraph{An exception: $\mathfrak{so}(8)$ as a de-facto continuum algebra from a discrete construction.}
Principled exception to the discrete-only stance: the closed Lie algebra $\mathfrak{so}(8) \supseteq G_2\cup\{\Delta\}$ generated from the Fano octonion left-multiplication table and the phase-lift generator $\Delta\in(\mathrm{e}_1,\mathrm{e}_7)$.  Although $\mathfrak{so}(8)$ is a continuous (28-dimensional real) Lie algebra, it is built here entirely from \emph{discrete} ingredients (the seven imaginary octonion units, a finite Fano table, and one $\pi/2$ rotation), and its closure to 28 dimensions is a finite linear-algebra certificate (see the \texttt{Hqiv.SO8Closure} / \texttt{Hqiv.SO8ClosureInterface} stack).  The Lie-group structure is therefore a \emph{de-facto continuum algebra obtained from a discrete construction}: continuous in parameter space, but with no continuum spacetime, no continuum measure, and no UV limit attached.  When this manuscript or its companions write $\exp(t\,\Delta)\in\mathrm{SO}(8)$ or treat $\mathfrak{so}(8)$ rotations as smooth, those are statements internal to a finite-dimensional Lie group whose generators are certified by discrete data; they are \emph{not} a continuum-spacetime commitment and do not require a continuum-QFT scaffolding.

% From papers/paper/*.tex: % Shared reader contract: patch QFT + discrete numerics.
% From papers/*.tex:     % Shared reader contract: patch QFT + discrete numerics.
% From papers/*.tex:     \input{include/patch_theory_messaging.tex}
% From papers/paper/*.tex: \input{../include/patch_theory_messaging.tex}

\paragraph{Patch QFT and discrete numerics (reader contract).}
HQIV (Horizon-Quantized Informational Vacuum) is formulated as a \emph{patch theory}: quantum fields, actions, and physical readouts are attached to discrete patches---shell indices $m\in\mathbb{N}$, finite spacetime index types (e.g.\ $\mathrm{Fin}\,4$), and horizon-limited charts---rather than to a hypothetical fundamental smooth spacetime obtained as a mesh-refinement or ultraviolet limit.
Accordingly, when this text refers to ``QFT,'' it means \emph{patch QFT}: patching rules, finite measurement ledgers, and variational identities on the discrete spine; it is not the claim that the theory is \emph{defined} by recovering ordinary continuum quantum field theory through such a limit.
The numbers that admit the sharpest statements---mass and coupling ladders, curvature ratios, shell thermodynamics, and rational witnesses in \texttt{HQIV\_LEAN}---are objects of \emph{discrete mathematics} on $\mathbb{N}$ (combinatorial growth, cumulative simplex ledgers) together with the ordered field $\mathbb{R}$ as used in the proof assistant for analysis, bounds, and phenomenological display.
Smooth-manifold or continuum field notation, when it appears, is a \emph{translation layer} for comparison with standard physics literature; it does not supersede the discrete definitions.

\paragraph{A continuum is not required, and in places would destroy these results.}
The discrete patch layer is \emph{not} a numerical approximation to a more fundamental continuum theory awaiting future construction; it is the ontology.  Where the literature treats ``the continuum limit'' as the goal, HQIV treats it as a \emph{calculation approximation} for comparison with standard formulas.  The continuum form is an optional convenience for comparison; the proved discrete statement is what carries the physics.  In several load-bearing places---the curvature imprint $\delta_E(m)$ at shell $m$, the harmonic ladder masses $m_\tau\to m_\mu\to m_e$, the lock-in vev $v=\sqrt{\eta_{\rm paper}\,\Omega_k(m_{\rm lockin};m_{\rm lockin})}$, the proton anchor at $m=\mathrm{referenceM}=4$, the baryon-asymmetry $\eta\propto 6^7\sqrt{3}/(m+1)$, and the rational coefficients $\alpha=3/5$, $\gamma=2/5$, $\alpha_{\rm GUT}=1/42$---taking a literal $m\to\infty$ or sub-Planck continuum limit would \emph{erase} the discrete index that is doing the predictive work.  Continuum forms of these objects are therefore not preferred refinements; they are degenerate, information-losing readouts of the same patch data.

\paragraph{An exception: $\mathfrak{so}(8)$ as a de-facto continuum algebra from a discrete construction.}
Principled exception to the discrete-only stance: the closed Lie algebra $\mathfrak{so}(8) \supseteq G_2\cup\{\Delta\}$ generated from the Fano octonion left-multiplication table and the phase-lift generator $\Delta\in(\mathrm{e}_1,\mathrm{e}_7)$.  Although $\mathfrak{so}(8)$ is a continuous (28-dimensional real) Lie algebra, it is built here entirely from \emph{discrete} ingredients (the seven imaginary octonion units, a finite Fano table, and one $\pi/2$ rotation), and its closure to 28 dimensions is a finite linear-algebra certificate (see the \texttt{Hqiv.SO8Closure} / \texttt{Hqiv.SO8ClosureInterface} stack).  The Lie-group structure is therefore a \emph{de-facto continuum algebra obtained from a discrete construction}: continuous in parameter space, but with no continuum spacetime, no continuum measure, and no UV limit attached.  When this manuscript or its companions write $\exp(t\,\Delta)\in\mathrm{SO}(8)$ or treat $\mathfrak{so}(8)$ rotations as smooth, those are statements internal to a finite-dimensional Lie group whose generators are certified by discrete data; they are \emph{not} a continuum-spacetime commitment and do not require a continuum-QFT scaffolding.

% From papers/paper/*.tex: % Shared reader contract: patch QFT + discrete numerics.
% From papers/*.tex:     \input{include/patch_theory_messaging.tex}
% From papers/paper/*.tex: \input{../include/patch_theory_messaging.tex}

\paragraph{Patch QFT and discrete numerics (reader contract).}
HQIV (Horizon-Quantized Informational Vacuum) is formulated as a \emph{patch theory}: quantum fields, actions, and physical readouts are attached to discrete patches---shell indices $m\in\mathbb{N}$, finite spacetime index types (e.g.\ $\mathrm{Fin}\,4$), and horizon-limited charts---rather than to a hypothetical fundamental smooth spacetime obtained as a mesh-refinement or ultraviolet limit.
Accordingly, when this text refers to ``QFT,'' it means \emph{patch QFT}: patching rules, finite measurement ledgers, and variational identities on the discrete spine; it is not the claim that the theory is \emph{defined} by recovering ordinary continuum quantum field theory through such a limit.
The numbers that admit the sharpest statements---mass and coupling ladders, curvature ratios, shell thermodynamics, and rational witnesses in \texttt{HQIV\_LEAN}---are objects of \emph{discrete mathematics} on $\mathbb{N}$ (combinatorial growth, cumulative simplex ledgers) together with the ordered field $\mathbb{R}$ as used in the proof assistant for analysis, bounds, and phenomenological display.
Smooth-manifold or continuum field notation, when it appears, is a \emph{translation layer} for comparison with standard physics literature; it does not supersede the discrete definitions.

\paragraph{A continuum is not required, and in places would destroy these results.}
The discrete patch layer is \emph{not} a numerical approximation to a more fundamental continuum theory awaiting future construction; it is the ontology.  Where the literature treats ``the continuum limit'' as the goal, HQIV treats it as a \emph{calculation approximation} for comparison with standard formulas.  The continuum form is an optional convenience for comparison; the proved discrete statement is what carries the physics.  In several load-bearing places---the curvature imprint $\delta_E(m)$ at shell $m$, the harmonic ladder masses $m_\tau\to m_\mu\to m_e$, the lock-in vev $v=\sqrt{\eta_{\rm paper}\,\Omega_k(m_{\rm lockin};m_{\rm lockin})}$, the proton anchor at $m=\mathrm{referenceM}=4$, the baryon-asymmetry $\eta\propto 6^7\sqrt{3}/(m+1)$, and the rational coefficients $\alpha=3/5$, $\gamma=2/5$, $\alpha_{\rm GUT}=1/42$---taking a literal $m\to\infty$ or sub-Planck continuum limit would \emph{erase} the discrete index that is doing the predictive work.  Continuum forms of these objects are therefore not preferred refinements; they are degenerate, information-losing readouts of the same patch data.

\paragraph{An exception: $\mathfrak{so}(8)$ as a de-facto continuum algebra from a discrete construction.}
Principled exception to the discrete-only stance: the closed Lie algebra $\mathfrak{so}(8) \supseteq G_2\cup\{\Delta\}$ generated from the Fano octonion left-multiplication table and the phase-lift generator $\Delta\in(\mathrm{e}_1,\mathrm{e}_7)$.  Although $\mathfrak{so}(8)$ is a continuous (28-dimensional real) Lie algebra, it is built here entirely from \emph{discrete} ingredients (the seven imaginary octonion units, a finite Fano table, and one $\pi/2$ rotation), and its closure to 28 dimensions is a finite linear-algebra certificate (see the \texttt{Hqiv.SO8Closure} / \texttt{Hqiv.SO8ClosureInterface} stack).  The Lie-group structure is therefore a \emph{de-facto continuum algebra obtained from a discrete construction}: continuous in parameter space, but with no continuum spacetime, no continuum measure, and no UV limit attached.  When this manuscript or its companions write $\exp(t\,\Delta)\in\mathrm{SO}(8)$ or treat $\mathfrak{so}(8)$ rotations as smooth, those are statements internal to a finite-dimensional Lie group whose generators are certified by discrete data; they are \emph{not} a continuum-spacetime commitment and do not require a continuum-QFT scaffolding.


\paragraph{Patch QFT and discrete numerics (reader contract).}
HQIV (Horizon-Quantized Informational Vacuum) is formulated as a \emph{patch theory}: quantum fields, actions, and physical readouts are attached to discrete patches---shell indices $m\in\mathbb{N}$, finite spacetime index types (e.g.\ $\mathrm{Fin}\,4$), and horizon-limited charts---rather than to a hypothetical fundamental smooth spacetime obtained as a mesh-refinement or ultraviolet limit.
Accordingly, when this text refers to ``QFT,'' it means \emph{patch QFT}: patching rules, finite measurement ledgers, and variational identities on the discrete spine; it is not the claim that the theory is \emph{defined} by recovering ordinary continuum quantum field theory through such a limit.
The numbers that admit the sharpest statements---mass and coupling ladders, curvature ratios, shell thermodynamics, and rational witnesses in \texttt{HQIV\_LEAN}---are objects of \emph{discrete mathematics} on $\mathbb{N}$ (combinatorial growth, cumulative simplex ledgers) together with the ordered field $\mathbb{R}$ as used in the proof assistant for analysis, bounds, and phenomenological display.
Smooth-manifold or continuum field notation, when it appears, is a \emph{translation layer} for comparison with standard physics literature; it does not supersede the discrete definitions.

\paragraph{A continuum is not required, and in places would destroy these results.}
The discrete patch layer is \emph{not} a numerical approximation to a more fundamental continuum theory awaiting future construction; it is the ontology.  Where the literature treats ``the continuum limit'' as the goal, HQIV treats it as a \emph{calculation approximation} for comparison with standard formulas.  The continuum form is an optional convenience for comparison; the proved discrete statement is what carries the physics.  In several load-bearing places---the curvature imprint $\delta_E(m)$ at shell $m$, the harmonic ladder masses $m_\tau\to m_\mu\to m_e$, the lock-in vev $v=\sqrt{\eta_{\rm paper}\,\Omega_k(m_{\rm lockin};m_{\rm lockin})}$, the proton anchor at $m=\mathrm{referenceM}=4$, the baryon-asymmetry $\eta\propto 6^7\sqrt{3}/(m+1)$, and the rational coefficients $\alpha=3/5$, $\gamma=2/5$, $\alpha_{\rm GUT}=1/42$---taking a literal $m\to\infty$ or sub-Planck continuum limit would \emph{erase} the discrete index that is doing the predictive work.  Continuum forms of these objects are therefore not preferred refinements; they are degenerate, information-losing readouts of the same patch data.

\paragraph{An exception: $\mathfrak{so}(8)$ as a de-facto continuum algebra from a discrete construction.}
Principled exception to the discrete-only stance: the closed Lie algebra $\mathfrak{so}(8) \supseteq G_2\cup\{\Delta\}$ generated from the Fano octonion left-multiplication table and the phase-lift generator $\Delta\in(\mathrm{e}_1,\mathrm{e}_7)$.  Although $\mathfrak{so}(8)$ is a continuous (28-dimensional real) Lie algebra, it is built here entirely from \emph{discrete} ingredients (the seven imaginary octonion units, a finite Fano table, and one $\pi/2$ rotation), and its closure to 28 dimensions is a finite linear-algebra certificate (see the \texttt{Hqiv.SO8Closure} / \texttt{Hqiv.SO8ClosureInterface} stack).  The Lie-group structure is therefore a \emph{de-facto continuum algebra obtained from a discrete construction}: continuous in parameter space, but with no continuum spacetime, no continuum measure, and no UV limit attached.  When this manuscript or its companions write $\exp(t\,\Delta)\in\mathrm{SO}(8)$ or treat $\mathfrak{so}(8)$ rotations as smooth, those are statements internal to a finite-dimensional Lie group whose generators are certified by discrete data; they are \emph{not} a continuum-spacetime commitment and do not require a continuum-QFT scaffolding.


\paragraph{Patch QFT and discrete numerics (reader contract).}
HQIV (Horizon-Quantized Informational Vacuum) is formulated as a \emph{patch theory}: quantum fields, actions, and physical readouts are attached to discrete patches---shell indices $m\in\mathbb{N}$, finite spacetime index types (e.g.\ $\mathrm{Fin}\,4$), and horizon-limited charts---rather than to a hypothetical fundamental smooth spacetime obtained as a mesh-refinement or ultraviolet limit.
Accordingly, when this text refers to ``QFT,'' it means \emph{patch QFT}: patching rules, finite measurement ledgers, and variational identities on the discrete spine; it is not the claim that the theory is \emph{defined} by recovering ordinary continuum quantum field theory through such a limit.
The numbers that admit the sharpest statements---mass and coupling ladders, curvature ratios, shell thermodynamics, and rational witnesses in \texttt{HQIV\_LEAN}---are objects of \emph{discrete mathematics} on $\mathbb{N}$ (combinatorial growth, cumulative simplex ledgers) together with the ordered field $\mathbb{R}$ as used in the proof assistant for analysis, bounds, and phenomenological display.
Smooth-manifold or continuum field notation, when it appears, is a \emph{translation layer} for comparison with standard physics literature; it does not supersede the discrete definitions.

\paragraph{A continuum is not required, and in places would destroy these results.}
The discrete patch layer is \emph{not} a numerical approximation to a more fundamental continuum theory awaiting future construction; it is the ontology.  Where the literature treats ``the continuum limit'' as the goal, HQIV treats it as a \emph{calculation approximation} for comparison with standard formulas.  The continuum form is an optional convenience for comparison; the proved discrete statement is what carries the physics.  In several load-bearing places---the curvature imprint $\delta_E(m)$ at shell $m$, the harmonic ladder masses $m_\tau\to m_\mu\to m_e$, the lock-in vev $v=\sqrt{\eta_{\rm paper}\,\Omega_k(m_{\rm lockin};m_{\rm lockin})}$, the proton anchor at $m=\mathrm{referenceM}=4$, the baryon-asymmetry $\eta\propto 6^7\sqrt{3}/(m+1)$, and the rational coefficients $\alpha=3/5$, $\gamma=2/5$, $\alpha_{\rm GUT}=1/42$---taking a literal $m\to\infty$ or sub-Planck continuum limit would \emph{erase} the discrete index that is doing the predictive work.  Continuum forms of these objects are therefore not preferred refinements; they are degenerate, information-losing readouts of the same patch data.

\paragraph{An exception: $\mathfrak{so}(8)$ as a de-facto continuum algebra from a discrete construction.}
Principled exception to the discrete-only stance: the closed Lie algebra $\mathfrak{so}(8) \supseteq G_2\cup\{\Delta\}$ generated from the Fano octonion left-multiplication table and the phase-lift generator $\Delta\in(\mathrm{e}_1,\mathrm{e}_7)$.  Although $\mathfrak{so}(8)$ is a continuous (28-dimensional real) Lie algebra, it is built here entirely from \emph{discrete} ingredients (the seven imaginary octonion units, a finite Fano table, and one $\pi/2$ rotation), and its closure to 28 dimensions is a finite linear-algebra certificate (see the \texttt{Hqiv.SO8Closure} / \texttt{Hqiv.SO8ClosureInterface} stack).  The Lie-group structure is therefore a \emph{de-facto continuum algebra obtained from a discrete construction}: continuous in parameter space, but with no continuum spacetime, no continuum measure, and no UV limit attached.  When this manuscript or its companions write $\exp(t\,\Delta)\in\mathrm{SO}(8)$ or treat $\mathfrak{so}(8)$ rotations as smooth, those are statements internal to a finite-dimensional Lie group whose generators are certified by discrete data; they are \emph{not} a continuum-spacetime commitment and do not require a continuum-QFT scaffolding.


\paragraph{Reproducibility archive.}
Unpack the release distribution \ReleaseDistArchive{} (frozen for \path{papers/closure.tex}). The Lean modules, scripts, JSON certificates, and PDF/\TeX{} sources cited below for audit are mirrored under \PaperReproBundleDir/ (see \PaperReproBundleManifest). Paths in this paper are relative to \PaperReproBundleDir/ after unpacking that subtree from the release archive.

\section{Introduction}

This paper presents one continuous geometric-to-algebraic derivation:
\[
\text{uniform growth} \Longrightarrow K \Longrightarrow \Omega \Longrightarrow \mathcal R
\Longrightarrow \Delta \Longrightarrow \langle G_2,\Delta\rangle_{\mathrm{Lie}}=\mathfrak{so}(8).
\]
Historically, the closure was first reached through the full rapidity-fiber construction; the cleaner \(G_2+\Delta\) closure theorem used here is the compressed algebraic form of that original route.

Write \emph{HQIV} for \emph{Horizon-Quantized Informational Vacuum}: the framework that ties entanglement monogamy on overlapping causal horizons to discrete null-shell combinatorics and related readouts, as developed in the companion physics manuscript and machine-checked in the \cite{hqiv-lean} library.
Within HQIV, the chosen growth law is physics-motivated by discrete mode counting for a local light-cone construction on flat background space.

\medskip
\noindent\textbf{Scope note (read first).}
\textbf{Proved in Lean with minimal hypotheses:} (i) positivity of \(\rho\) and divergence of \(K\) for every fixed \(\alpha>0\) (harmonic lower bound \(K(n)\ge H_n\)); (ii) holonomy identities for embedded plane generators; (iii) Lie closure of the \emph{explicit} \(G_2\cup\{\Delta\}\) seed to \(\mathfrak{so}(8)\).
\textbf{Imported for motivation (not a consequence of poset axioms):} the quadratic law \(A(m)=4(m+2)(m+1)\), the value \(\alpha=3/5\), the Fano-basis octonion model, the color axis \(e_7\), the \((e_1,e_7)\) phase plane, and the Brodie thermodynamic story \cite{brodie2026}.
A one-paragraph HQIV synopsis and a Lean dependency sketch appear in Appendix~\ref{app:hqiv-box}.

The same analytic mechanism applies to any growth law that yields a positive divergent cumulative channel, including non-physics settings.\,\footnote{In particular, Theorem~\ref{thm:main-curvature}'s harmonic lower bound and divergence use only \(\alpha>0\) and the positivity of \(\rho\); replacing the HQIV quadratic ladder by a different discrete law would change the numerical details of \(K\) and \(\Omega\) while preserving the same \emph{template} whenever the new law still yields a divergent normalized channel.}

\begin{remark}[Physics motivation vs.\ algebraic independence]
The HQIV narrative and companion manuscript motivate the \emph{choice} of ladder and normalization; they are not logical prerequisites for the Lie-closure certificate once the phase-lift generator \(\Delta\) and octonionic seed are fixed. Concretely: the machine-checked statement \(G_2+\Delta\) closes to \(\mathfrak{so}(8)\) is a theorem of linear algebra and Lie brackets over explicit matrices; the curvature/rapidity chain is a separate geometric template that supplies a natural route to \(\Delta\) in the HQIV setting. This paper does \emph{not} claim that discrete shell growth physically \emph{causes} \(\mathfrak{so}(8)\) to appear in nature---only that the same minimal-input \(\Rightarrow\) rigid-output pattern appears in both the growth-to-rapidity layer and the seed-to-closure layer (see Proposition~\ref{thm:forced-structure} and the remark following it).
\end{remark}

Let $m\in\mathbb N$ be a discrete index parameter and define
\[
N(m):=(m+2)(m+1),\qquad A(m):=4N(m),
\]
with fixed exponent
\[
\alpha:=\frac35.
\]
We explicitly credit this \(\alpha=3/5\) normalization to Brodie's thermodynamic horizon analysis \cite{brodie2026}; the lattice-ratio identity is the formal verification anchor in the companion Lean development.
Readers should treat the Zenodo record as the authoritative statement of Brodie's derivation and verify the DOI at submission time.

\medskip
\noindent\textbf{Robustness (central).} For the analytic growth\(\to\)curvature chain \emph{alone}, any fixed \(\alpha>0\) preserves positivity of \(\rho\) on \(x\ge 1\), the harmonic lower bound \(K(n)\ge H_n\), and divergence \(K(n)\to\infty\); the value \(\alpha=3/5\) is the HQIV-distinguished normalization, not a logical necessity for Theorem~\ref{thm:main-curvature}.
Define curvature density and cumulative curvature
\[
\rho(x):=\frac1x\bigl(1+\alpha\log x\bigr)\quad (x\ge 1),\qquad
K(n):=\sum_{m=0}^{n-1}\rho(m+1).
\]
Fix any normalization index $m_\ast$ with $K(m_\ast)>0$ and set
\[
\Omega(n):=\frac{K(n)}{K(m_\ast)}.
\]
Any index \(m_\ast\) with \(K(m_\ast)>0\) is admissible; all identities below are independent of which admissible choice is made. In the companion Lean library, normalization is pinned at \texttt{referenceM}\(=4\) (a calibration convention), but the proofs use only positivity of \(K(\texttt{referenceM})\).

\subsection*{Quadratic null-shell counting versus cubic CST ball growth}

In Lorentzian dimension \(3+1\), a causal diamond whose proper height is \(T\) has spatial volume of order \(T^3\), so counting elements in coarse-grained shells often scales as a \emph{cubic} polynomial in the continuum limit. The HQIV null-lattice law used here is instead the fixed \emph{quadratic} increment law of Proposition~\ref{prop:growth-identities},
\[
A(m)=4(m+2)(m+1),\qquad A(m+1)-A(m)=8(m+2),
\]
motivated by stars-and-bars counting on a discrete light cone (companion HQIV manuscript). The present paper should therefore be read as conditioning on that HQIV discretization choice: replacing \(A(m)\) by a different polynomial (for example, one tracking bulk ball volume rather than null-shell mode unlocks) changes the constants in \(\rho\) and \(K\) and may require revisiting the normalization story, even though the \(\alpha>0\) harmonic/divergence skeleton persists.

\subsection*{Historical derivation arc}

The development history matters for interpretation. The original route was
\[
\text{rapidity fiber} \;\longrightarrow\; \text{phase-lift direction } \Delta \;\longrightarrow\; \mathfrak{so}(8)\text{ closure}.
\]
In that route, the rapidity phase channel provides the geometric phase direction whose infinitesimal action is encoded by \(\Delta\). The present paper then uses a cleaner algebraic route: once \(\Delta\) is isolated, adjoining it to the 14 generators of \(G_2\) yields the full closure statement.

This should be read as a compression, not a replacement: the concise \(G_2+\Delta\) theorem is the distilled algebraic core of the earlier phase-fiber derivation. The formal claims in this manuscript concern the compressed route; the historical geometric motivation is included to keep the two parts of the story logically connected.

\subsection*{Causal Set Theory as Discrete Substrate for Emergent Gauge Structure}

Causal set theory (CST) furnishes a fundamentally discrete model of spacetime: a locally finite partially ordered set (poset) whose continuum limit, under suitable sprinkling, recovers a Lorentzian manifold. In CST the fundamental datum is the causal relation; derived quantities include the cardinality of causal intervals and the growth of ``shells''---the number of elements at successive discrete proper-time steps from a reference event. For a 3+1-dimensional causal set approximating flat or curved spacetime, this shell-counting function \(A(m)\) is polynomial of degree 3 in the continuum limit (volume \(\sim t^3\)), but the effective null-shell or light-cone discretization relevant to HQIV reduces to the quadratic growth law adopted in this paper:
\[
A(m)=4(m+2)(m+1).
\]
Any discrete model whose \emph{shell increment} matches the hypothesis of Section~\ref{sec:growth} (in particular, the same quadratic \(A(m)\) and a choice of \(\alpha>0\) in \(\rho\)) is compatible with the analytic chain \(K,\Omega,\mathcal{R}\) of Sections~\ref{sec:growth}--\ref{sec:rapidity}. In particular:
\begin{itemize}
  \item such a law yields a positive divergent cumulative curvature channel \(K\) (Theorem~\ref{thm:main-curvature});
  \item normalization produces the normalized cumulative phase readout \(\mathcal{R}\) once phase-map axioms are fixed;
  \item a phase increment defines a phase-lift generator \(\Delta\) only after one commits to an explicit \(2\)-plane in \(\mathbb{R}^8\) (here the HQVM \((e_1,e_7)\) choice);
  \item given the \emph{same imported} octonionic matrices \(G_2\cup\{\Delta\}\) used in Section~\ref{sec:closure}, Lie closure to \(\mathfrak{so}(8)\) is a certified algebraic fact.
\end{itemize}
Causal set theory does \emph{not} by itself determine that multiplication table, axis, or plane; poset combinatorics alone does not replace those inputs. What \emph{is} rigid, once those inputs are fixed, is the template ``divergent \(K\) \(\Rightarrow\) normalized phase readout \(\Rightarrow\) planar holonomy \(\Rightarrow\) Lie closure of the chosen seed.'' The Fourier-twiddle construction on \(S^{n-1}\) extends the holonomy layer to general \(\mathrm{SO}(n)\) once a seed subalgebra is chosen there as well.

\paragraph{Generic minimal \(\mathfrak{so}(N)\) seed (coordinate model).}\label{para:minimal-so-seed}
Fix \(N\ge 3\). Embed \(\mathfrak{so}(N-1)\) as skew-adjoint matrices supported on the first \(N-1\) coordinates, and adjoin one plane generator \(\Delta_{k,N}=E_{k,N}-E_{N,k}\) connecting a chosen index \(k\) in that block to the final basis vector (the \(N=4\) quaternionic toy model in Section~\ref{sec:closure} is this pattern for the standard labeling). The Lie subalgebra generated by this finite seed equals the full matrix Lie algebra \(\mathfrak{so}(N)\) in the standard skew-adjoint model on \(\mathbb{R}^N\). A machine-checked proof (for all \(N\ge 3\) and every admissible \(k\)) is in the HQIV Lean library: \path{Hqiv/Algebra/MinimalSoSeedClosure.lean}, main statement \texttt{minimal\_so\_seed\_lieSpan\_eq\_skewAdjoint}. For \(N=8\) one may still use the smaller exceptional \(G_2+\Delta\) seed recorded elsewhere in the library; the coordinate seed above is the dimension-agnostic default aligned with the Fourier-twiddle plane family.

\section{Uniform Growth and the Forced Curvature Channel}
\label{sec:growth}
This is Step 1 of the derivation: growth law \(\to \rho(x)\to K(n)\to \Omega(n)\).

\begin{proposition}[Exact growth identities]
\label{prop:growth-identities}
For all $m\in\mathbb N$,
\[
A(m)=4m^2+12m+8,\qquad A(m+1)-A(m)=8(m+2).
\]
\end{proposition}

\begin{proof}
Expand $A(m)=4(m+2)(m+1)$ and subtract successive values.
\end{proof}

\begin{remark}[Why this $\rho$ and why $\alpha=\frac35$]
\label{rem:rho-motivation}
Proposition~\ref{prop:growth-identities} shows that the discrete derivative of the growth law is linear:
$A(m+1)-A(m)=8(m+2)$. A scale-invariant cumulative channel compatible with linear step growth is naturally of logarithmic-harmonic form
\[
\rho(x)=\frac{c+d\log x}{x},
\]
whose partial sums produce harmonic and log-corrected terms. The normalization used in this paper is $c=1$ and $d=\alpha=\frac35$, following Brodie's horizon-thermodynamic coefficient choice \cite{brodie2026}, matching the exact lattice-ratio identity, and giving immediate positivity on $x\ge 1$.
Within the HQIV framework this quadratic ladder is the discrete light-cone shell model; outside that framework it can be replaced by alternative growth laws while preserving the same proof pattern whenever positivity/divergence hypotheses are retained.
\end{remark}

\begin{lemma}[Positivity of density]
For all $x\ge 1$, $\rho(x)>0$.
\end{lemma}

\begin{proof}
Since $x\ge 1$, $\log x\ge 0$. Therefore $1+\alpha\log x\ge 1>0$ and $1/x>0$, so $\rho(x)>0$.
\end{proof}

\begin{remark}[Parameter robustness]
The proofs of positivity of \(\rho\), the lower bound \(K(n)\ge H_n\), and therefore divergence of \(K\), require only \(\alpha>0\) (indeed \(\alpha\ge 0\) suffices for the lower bound). Thus the logical chain
\[
\text{growth}\to \rho \to K \to \Omega \to \mathcal R
\]
is structurally robust under positive choices of \(\alpha\); the value \(3/5\) is a model-specific normalization rather than a necessity for the abstract analytic mechanism.
\end{remark}

\begin{theorem}[Uniform-growth curvature theorem]
\label{thm:main-curvature}
For all $n\ge 1$,
\[
K(n)\ge H_n:=\sum_{m=0}^{n-1}\frac1{m+1}.
\]
Hence $K(n)\to +\infty$ as $n\to\infty$.
\end{theorem}

\begin{proof}
By definition,
\[
\rho(m+1)=\frac1{m+1}\bigl(1+\alpha\log(m+1)\bigr)\ge \frac1{m+1}
\]
because $\log(m+1)\ge 0$. Combined with the linear step-growth from Proposition~\ref{prop:growth-identities} and the logarithmic-harmonic ansatz in Remark~\ref{rem:rho-motivation}, summing gives $K(n)\ge H_n$.
Since $H_n$ diverges, so does $K(n)$.
\end{proof}

\begin{remark}
The conclusion is pointwise and asymptotic: each finite partial sum is bounded below by harmonic growth, and the entire channel diverges.
\end{remark}

\section{Rapidity as a Primitive Phase Observable}
\label{sec:rapidity}
This is Step 2: \(\Omega\to\mathcal R\). The formal object is a \emph{normalized cumulative phase readout} built from \(\Omega\) and a phase map \(\mathcal F\) satisfying the axioms below; we refer to it as ``rapidity'' only to flag the interpretive connection to classical \(\operatorname{artanh}v\) and boost parametrizations. No claim is made that CST axioms or discrete posets uniquely identify \(\mathcal{F}\) with relativistic rapidity, nor that Berry/Aharonov--Anandan phases are recovered without further structure.

\begin{definition}[Normalized curvature readout]
Let $\Omega(n):=K(n)/K(m_\ast)$ with $K(m_\ast)>0$.
\end{definition}

\begin{lemma}
\[
\Omega(m_\ast)=1,\qquad \Omega(n)\to +\infty \text{ as } n\to\infty.
\]
\end{lemma}

\begin{proof}
The first identity is immediate from normalization. The second follows from
Theorem~\ref{thm:main-curvature} and positivity of $K(m_\ast)$.
\end{proof}

\begin{definition}[Rapidity phase channel]
Define the \emph{normalized cumulative phase readout} (``rapidity'' in interpretive language) by
\[
\mathcal R(\phi,t,n):=\mathcal F(\phi,t,\Omega(n)),
\]
where $\mathcal F$ is the phase map. At $n=m_\ast$, this recovers the base harmonic phase.
\end{definition}

The readout \(\mathcal R\) is determined, at the level of proofs in this manuscript, by the axioms on \(\mathcal F\) together with divergence of \(\Omega\); additional physics would be needed to justify the name ``rapidity'' beyond analogy.

\begin{remark}[Phase-map axioms: minimal interface]
\label{rem:phase-map-interface}
The following three bullets are the \emph{minimal interface} assumed in Sections~\ref{sec:growth}--\ref{sec:closure}: any concrete realization (Remark~\ref{rem:phase-map-concrete}) must imply them, but the proofs never use a closed form beyond these properties.
For the derivations below, $\mathcal F$ is only required to satisfy:
\begin{itemize}
  \item normalization at the reference index: $\mathcal F(\phi,t,1)=\mathcal H(\phi,t)$;
  \item monotonic dependence on curvature channel: $\Omega_1\le \Omega_2 \Rightarrow \mathcal F(\phi,t,\Omega_1)\le \mathcal F(\phi,t,\Omega_2)$ (for fixed $(\phi,t)$);
  \item embedding invariance: for any embedding $e$, pullback commutes with evaluation, i.e. $(\mathcal F\circ(\phi,t,\Omega))\circ e=\mathcal F\circ(\phi\circ e,t\circ e,\Omega\circ e)$.
\end{itemize}
No stronger closed form is needed for the closure and transport statements in this manuscript.
If \(\mathcal F\) is not monotone in \(\Omega\) (or is discontinuous in the channel variable), the monotonicity axiom above fails; comparison statements that implicitly order two channels would then require different hypotheses (e.g.\ Lipschitz or BV regularity on compact \(\Omega\)-intervals). The proofs in this manuscript assume exactly the axioms listed.
Continuum limits, smooth bundles, and smooth time parameters are \emph{interpretive}: they are not used anywhere in Sections~\ref{sec:growth}--\ref{sec:closure}; the formal chain is discrete in \(n\in\mathbb N\) throughout.
\end{remark}

\begin{remark}[Concrete phase-map instance]
\label{rem:phase-map-concrete}
To avoid black-box ambiguity, one admissible concrete choice is
\[
\mathcal F(\phi,t,\Omega):=\Omega\cdot \mathcal H(\phi,t),
\]
with \(\mathcal H\) the base harmonic phase. This satisfies normalization
\(\mathcal F(\phi,t,1)=\mathcal H(\phi,t)\), is monotone in \(\Omega\) whenever
\(\mathcal H(\phi,t)\ge 0\), and realizes the same discrete chain
\(\Omega(n)\mapsto \mathcal R(\phi,t,n)\).
\end{remark}

\begin{remark}[Interpretive bridge to classical rapidity and geometric phase]
The object \(\mathcal R(\phi,t,n)\) is discrete in this manuscript. It can be \emph{compared} to a hyperbolic-angle parameter (classical relativistic rapidity) because both organize one-parameter families of phase-like updates; likewise, monodromy in \(\Omega\) suggests loose analogies with Berry/Aharonov--Anandan phases. These parallels are \emph{not} proved here as equivalences: CST and thermodynamic inputs would be required to promote them from metaphor to theorem.
\end{remark}

\paragraph{Interpretation vs.\ proofs (bridge language).}
No continuum limit, Lorentzian metric, or relativistic boost interpretation is assumed or proved for \(\mathcal R\): the formal chain is indexed by \(n\in\mathbb N\) and uses only the phase-map axioms together with positivity and divergence of \(\Omega\). The vocabulary ``primitive rapidity'' means ``first readout after normalization of \(K\),'' not ``unique classical rapidity coordinate in a given spacetime.'' The same clarification applies in Section~\ref{sec:cst-forcing}: causal-set language motivates indexing, not a continuum field theory of \(\mathcal R\).

\begin{remark}
The construction remains explicitly chained
\[
\text{growth}\to K\to \Omega\to \mathcal R,
\]
but the role of \(\mathcal R\) is primitive at the level of output type: once \(\mathcal F\) is fixed, it is the first-order phase readout attached to any divergent normalized channel in this framework.
\end{remark}

\subsection*{Sign-flip invariance}
The phase readout built from \((\rho,K,\Omega)\) is compatible with the following discrete symmetries of the data:
\[
\rho\mapsto -\rho,\quad K\mapsto -K,\quad \Omega\mapsto \pm \Omega,\quad \theta\mapsto -\theta.
\]
At generator level, \(\Delta\mapsto -\Delta\) changes orientation of the same distinguished plane action and preserves \(2\pi\)-periodic closure.
\textbf{Index/reflection pairing \([x,|x|]\) (one-line convention).}
In this manuscript the bracket notation \([x,|x|]\) is shorthand for the \emph{index-growth reflection bookkeeping} compatible with the sign-flip package above: one pairs a signed channel coordinate \(x\) (e.g.\ a logarithmic or reciprocal readout built from shell labels) with its magnitude twin \(|x|\) when enforcing reversal symmetries on \((\rho,K,\Omega,\theta)\). No separate analytic formula for \([x,|x|]\) beyond this discrete compatibility is used in the proofs; continuum or critical-line realizations are left to companion work.

\subsection*{Universal outputs in the present framework}
Given a divergent cumulative channel and the phase-lift construction of Section~4, the framework-level outputs are:
\begin{itemize}
  \item exact \(2\pi\)-periodic plane holonomy generated by \(\Delta\);
  \item reflection pairing symmetry in index-growth form (\([x,|x|]\) / reciprocal pairing language);
  \item full Lie-closure to \(\mathfrak{so}(8)\) (or more generally to \(\mathfrak{so}(n)\)) once a suitable seed subalgebra is adjoined to the phase-lift generator (see Theorem~\ref{thm:universal-holonomy} and Section~5).
\end{itemize}

\begin{remark}[Discrete channel and ``fiber'' language]
The core proofs in this paper are discrete (\(n\in\mathbb N\)). The term ``rapidity fiber'' is used in the minimal sense of a one-parameter phase direction attached to each index state through \(\mathcal R(\phi,t,n)\), not as a fully developed smooth bundle construction. A continuum/fiber-bundle realization is treated as an interpretation layer motivated by embedding invariance and thermodynamic-limit heuristics, not as an additional theorem used in Sections~2--5.
\end{remark}

An optional commutative-diagram transport interface (not used in the discrete core) is recorded in Appendix~\ref{app:transport}.

\subsection*{Generalization to \texorpdfstring{$\mathrm{SO}(n)$}{SO(n)} via Fourier twiddles}
The phase-lift mechanism is dimension-independent once a distinguished \(2\)-plane is fixed. For \(n\ge 2\), choose \(i<j\) and let \(\Delta_{ij}\in\mathfrak{so}(n)\) be the standard antisymmetric plane-rotation generator (nonzero \(2\times2\) block on \(\mathrm{span}\{e_i,e_j\}\), zero elsewhere). For rapidity increment \(\theta(n)\), define
\[
U(\varepsilon,n):=\exp\!\bigl(\varepsilon\,\theta(n)\,\Delta_{ij}\bigr)\in \mathrm{SO}(n).
\]
On \(L^2(S^{n-1})\), the induced weight-\(k\) multiplier is the twiddle character
\[
\chi_k(\theta(n))=\exp\!\bigl(i\,k\,\theta(n)\bigr),\qquad k\in\mathbb Z.
\]

\medskip
\noindent\textbf{Caveat (holonomy vs.\ full closure).} Theorem~\ref{thm:universal-holonomy} below is a planar holonomy statement for an embedded generator \(\Delta_{ij}\in\mathfrak{so}(n)\). For \(n>2\), the Lie subalgebra generated by \(\{\Delta_{ij}\}\) alone is one-dimensional (the abelian span of \(\Delta_{ij}\)); bracket closure to the \emph{entire} \(\mathfrak{so}(n)\) still requires a suitable seed subalgebra and a certificate. Paragraph~\ref{para:minimal-so-seed} records one dimension-agnostic seed together with a Lean proof for all \(n\ge 3\). The \(G_2\cup\{\Delta\}\) package of Section~\ref{sec:closure} is the smaller exceptional alternative at \(n=8\) and should not be read as a substitute for generic-seed data when \(n\neq 8\).

\begin{theorem}[Universal phase-lift holonomy equation]
\label{thm:universal-holonomy}
For every \(n\ge 2\) and every embedded plane generator \(\Delta_{ij}\),
\[
\exp(2\pi\,\Delta_{ij})=I_n,
\]
equivalently
\[
\operatorname{tr}\!\bigl(\exp(t\Delta_{ij})\bigr)=n-2+2\cos t,
\]
so \(\operatorname{tr}(\exp(t\Delta_{ij}))=n\) whenever \(t=2\pi k\), \(k\in\mathbb Z\).
\end{theorem}

\begin{proof}
\(\Delta_{ij}\) is conjugate inside \(\mathfrak{so}(n)\) to the block matrix with
\(\begin{psmallmatrix}0&1\\-1&0\end{psmallmatrix}\) on one \(2\)-plane and zeros elsewhere. Exponentiation gives one planar rotation block \(\begin{psmallmatrix}\cos t&\sin t\\-\sin t&\cos t\end{psmallmatrix}\) and identity on the orthogonal complement, yielding the trace formula. At \(t=2\pi k\), the rotation block is identity, so the full matrix is \(I_n\).
The trace identity recovers the classical \(2\pi\)-periodicity of rotations in any embedded plane.
\end{proof}

\begin{remark}
The \(\mathfrak{so}(8)\) case is distinguished by the octonionic seed choice (\(G_2\cup\{\Delta\}\)); the twiddle/holonomy mechanism itself is canonical for all \(n\ge 2\) in this framework.
\end{remark}

\section{The Phase-Lift Generator \texorpdfstring{$\Delta$}{Delta}}

This is the bridge step from geometry to algebra. In the historical route, rapidity is first constructed as a phase observable \(\mathcal R\), and its infinitesimal phase direction is represented by the antisymmetric generator \(\Delta\). The concise route used for closure starts directly from this same \(\Delta\).

\subsection*{Explicit construction from the rapidity channel}
Fix a reference index \(m_\ast\) and define the rapidity phase increment
\[
\theta(n):=\mathcal R(\phi,t,n)-\mathcal R(\phi,t,m_\ast).
\]
On the distinguished phase plane \(P=\mathrm{span}\{e_1,e_7\}\subset\mathbb R^8\), define the one-parameter action
\[
U(\varepsilon,n):=\exp\!\bigl(\varepsilon\,\theta(n)\,\Delta\bigr),\qquad U(0,n)=I_8.
\]
The infinitesimal generator is
\[
\left.\frac{d}{d\varepsilon}U(\varepsilon,n)\right|_{\varepsilon=0}
=\theta(n)\,\Delta.
\]
In the standard basis, the phase-lift generator used in the closure seed is the skew-adjoint matrix with only two nonzero entries:
\[
\Delta_{1,7}=1,\qquad \Delta_{7,1}=-1,\qquad \Delta_{ij}=0\ \text{otherwise},
\]
equivalently the \(2\times2\) rotation block
\(
\begin{psmallmatrix}0&1\\-1&0\end{psmallmatrix}
\)
on \(P\) and zero on \(P^\perp\). Thus \(\Delta\in\mathfrak{so}(8)\) is the infinitesimal phase-direction lift of the rapidity channel.

\begin{remark}[Why the \(e_1\)-\(e_7\) plane]
The \(e_1\)-\(e_7\) phase plane is the canonical choice used by \texttt{HQVM/matrices.py}: \(e_7\) is the color-preferred octonionic axis in the fixed Fano-basis realization, and \(\Delta\) is defined as the infinitesimal rotation mixing \(e_1\) with \(e_7\). This aligns the geometric phase lift with the same basis conventions used by the \(G_2\) generator stack and triality-aware embedding data.
In the same multiplication-table convention, \(e_7\) is the \(G_2\)-stabilized color axis while \(e_1\) is the orthogonal phase/electromagnetic direction, so the \(e_1\)-\(e_7\) plane is the minimal triality-compatible phase-lift plane in this basis.
\end{remark}

\begin{remark}[Algebraic role of \(\Delta\)]
Beyond antisymmetry, the closure mechanism requires that \(\Delta\) contributes a direction not exhausted by the raw linear span of \(G_2\). This is certified in the formal development by \path{finrank_span_G2_union_Delta_le_15} together with \path{exists_so8Generator_not_mem_span_G2_union_Delta}, and completed by \path{lieBracket_in_span} and \path{G2_plus_Delta_closes_to_so8}.
\end{remark}

\begin{proposition}[Generator-level equivalence bridge]
The historical geometric route and the concise algebraic route are equivalent at generator level:
\[
\bigl(\text{rapidity fiber}\Rightarrow \Delta\bigr)
\Longleftrightarrow
\bigl(\Delta\ \text{used as the algebraic seed in }G_2+\Delta\bigr).
\]
\end{proposition}

\begin{proof}
Both routes identify the same phase-lift direction \(\Delta\). The geometric route derives it from the rapidity phase channel; the algebraic route takes that identified \(\Delta\) as input to Lie closure. Thus they differ by presentation length, not by generator content.
\end{proof}

\section{Algebraic Closure: \texorpdfstring{$G_2+\Delta\to\mathfrak{so}(8)$}{G2 + Delta to so(8)}}
\label{sec:closure}
This is Step 3 and the algebraic centerpiece: the distilled closure theorem. For a minimal standalone PDF collecting the symbolic generator scripts and JSON excerpts, see \path{papers/so8_closure_full_appendix.tex} (compiled to \path{papers/so8_closure_full_appendix.pdf}), included in \PaperArtifactBundleZip{} frozen for \path{papers/closure.tex}.

\subsection{Algebraic frame}
Let $\mathfrak g_2$ denote the 14-generator octonionic derivation sector and $\Delta$ the phase-lift antisymmetric generator in $M_8(\mathbb R)$.
Let $\{X_k\}_{k=1}^{28}$ denote the explicit antisymmetric basis candidate.

\begin{lemma}[Ambient model]
The space of real skew-adjoint $8\times8$ matrices has dimension $28$:
\[
\dim\{M\in M_8(\mathbb R):M^{\mathsf T}=-M\}=28.
\]
\end{lemma}

\begin{proof}
A skew-adjoint matrix is determined by strict upper-triangle entries; there are $\binom82=28$ such coordinates.
\end{proof}

\begin{lemma}[Seed containment]
Every generator in $\mathfrak g_2$ and $\Delta$ lies in the skew-adjoint model, hence in $\mathfrak{so}(8)$.
\end{lemma}

\begin{proof}
Each seed satisfies $M+M^{\mathsf T}=0$ by construction.
\end{proof}

\begin{lemma}[Linear-span obstruction]
The linear span of the raw seed set $G_2\cup\{\Delta\}$ has dimension at most $15$, so it cannot equal $\mathfrak{so}(8)$.
\end{lemma}

\begin{proof}
At most 15 listed matrices are linearly combined directly; but $\dim\mathfrak{so}(8)=28$.
\end{proof}

\begin{theorem}[Full Lie closure]
\label{thm:full-lie-closure}
The Lie algebra generated by $G_2\cup\{\Delta\}$ is all of $\mathfrak{so}(8)$.
\end{theorem}

\begin{proof}
Let $\mathfrak h:=\langle G_2,\Delta\rangle_{\mathrm{Lie}}$.

\emph{Step 1.} The candidate family $\{X_k\}_{k=1}^{28}$ is linearly independent and closed under commutator:
\[
[X_i,X_j]=\sum_{k=1}^{28} c_{ij}^k X_k.
\]
Hence $\mathrm{span}\{X_k\}$ is a 28-dimensional Lie subalgebra of skew-adjoint matrices.

\emph{Step 2.} Iterated commutators from $G_2\cup\{\Delta\}$ generate all $28$ basis directions (closure certificate), so $\dim\mathfrak h=28$.

\emph{Step 3.} Since $\mathfrak h\subseteq\mathfrak{so}(8)$ and both have dimension $28$, equality follows:
\[
\mathfrak h=\mathfrak{so}(8).
\]
\end{proof}

\begin{remark}
The key distinction is structural: the \emph{linear} seed span is not full, but the \emph{Lie-generated} span is full.
\end{remark}

\subsection*{Toy model: quaternionic closure \texorpdfstring{$\mathfrak{so}(3)+\Delta_4\to\mathfrak{so}(4)$}{so(3)+Delta4 to so(4)}}
The same linear-vs-Lie mechanism appears in the smallest nontrivial setting. Let
\[
J_{12},\ J_{13},\ J_{23}\in\mathfrak{so}(4)
\]
be the standard \(\mathfrak{so}(3)\) seed acting on the first three coordinates, and let
\[
\Delta_4:=J_{14}
\]
be the phase-lift generator (direct analog of the distinguished phase-plane lift in Section~4). Then:
\begin{itemize}
  \item the raw linear span of \(\{J_{12},J_{13},J_{23},\Delta_4\}\) has dimension \(4<6=\dim\mathfrak{so}(4)\);
  \item iterated brackets generate the missing directions \(J_{24},J_{34}\), since
  \[
  [J_{12},J_{14}]=J_{24},\qquad [J_{13},J_{14}]=J_{34};
  \]
  \item therefore \(\langle \mathfrak{so}(3),\Delta_4\rangle_{\mathrm{Lie}}=\mathfrak{so}(4)\).
\end{itemize}
So the closure pattern is already visible in a hand-checkable \(6\)-dimensional case before the octonionic \(28\)-dimensional setting (the same linear-vs-Lie mechanism and \(2\pi\)-holonomy hold verbatim for any \(\mathrm{SO}(n)\)).

In the octonionic certificate, the same phenomenon is visible in a single Lie word: for the machine-numbered basis \(\{X_k\}_{k=0}^{27}\) used in Lean, the bracket \([X_0,X_1]\) expands with nonzero coefficient on \(X_{27}\), e.g.
\[
c_{01}^{27}=\frac{-1315768990787}{25000000000000}\neq 0,
\]
so \([X_0,X_1]\notin \mathrm{span}(G_2\cup\{\Delta\})\) even though \([X_0,X_1]\) lies in the Lie closure (this is the same ``outside the raw seed span'' mechanism as \([J_{12},J_{14}]=J_{24}\), but now certified numerically inside the \(28\)-dimensional basis).

For reproducibility, the same exact-rational symbolic pipeline used for the \(\mathfrak{so}(8)\) certificate is mirrored in a toy script (\path{scripts/generate_symbolic_so4_certificate.py}) writing \path{artifacts/so4_symbolic_certificate.json} (paths relative to \PaperArtifactBundleDir/; archive \PaperArtifactBundleZip{} frozen for \path{papers/closure.tex}).

\subsection*{Explicit bracket strategy (illustrative fragment)}
The closure certificate is built by working in a fixed skew-adjoint basis
$\{X_k\}_{k=1}^{28}$ and computing all commutators in coordinates:
\[
[X_i,X_j]=\sum_{k=1}^{28} c_{ij}^k X_k.
\]
The computational pipeline constructs the $28\times 28$ basis matrix, inverts it once, and then projects each bracket vector to recover the coefficients $c_{ij}^k$. Lean then checks the resulting coefficient table inside the formal algebra.

A representative fully explicit row (from the generated closure table) is:
\[
[X_0,X_1]
=\sum_{k=0}^{27} c_{01}^k X_k,
\]
with sample coefficients
\[
c_{01}^2=\frac{-13165551793093}{100000000000000},\quad
c_{01}^{18}=\frac{-17776643565617}{200000000000000},\quad
c_{01}^{27}=\frac{-1315768990787}{25000000000000},
\]
and exact reconstruction identity in the symbolic certificate
\[
[X_0,X_1]-\sum_k c_{01}^kX_k
=0.
\]
This demonstrates the certificate pattern used for all pairs $(i,j)$.

For searchability in rendered PDFs, we also record the exact textual excerpt
emitted by the generator script:
\begin{quote}\ttfamily\small
| 0, 1, 2 => ((-13165551793093 : R) / 100000000000000)\\
| 0, 1, 18 => ((-17776643565617 : R) / 200000000000000)\\
| 0, 1, 27 => ((-1315768990787 : R) / 25000000000000)\\
-- Verification [0,1]:\\
max |bracket - sum c\_k gen\_k| = 2.220446049250313e-16
\end{quote}
These lines are produced by \path{scripts/print_lie_bracket_closure.py}
(\texttt{main()}, bracket check near the final verification print block).
Reproducibility note: this Python script is a coefficient generator (not itself
formally verified). Formal verification occurs in Lean by importing the emitted
tables (\path{Hqiv/GeneratorsLieClosureData0.lean}--\path{Hqiv/GeneratorsLieClosureData27.lean},
assembled by \path{Hqiv/GeneratorsLieClosureData.lean}) and proving
\path{lieBracket_in_span} and \path{G2_plus_Delta_closes_to_so8}.
The coefficient table is split across twenty-eight small Lean modules \path{Hqiv/GeneratorsLieClosureData0.lean}, \ldots, \path{Hqiv/GeneratorsLieClosureData27.lean} purely as an engineering measure (elaboration stack and tooling limits on very large proof objects); the mathematics is unchanged from a monolithic table.
Thus the certificate is machine-checked at theorem level even though coefficient
production is computational.
In parallel, we now provide an exact-arithmetic symbolic proof-object generated
by \path{scripts/generate_symbolic_lie_closure_certificate.py}, exported as
\path{artifacts/so8_symbolic_certificate.json}. That artifact records a rank-28
basis and all structure coefficients over \(\mathbb Q\) (numerator/denominator
pairs), with exact linear solves for each bracket relation.

\begin{remark}[Methodological caveat: numerical certificate vs exact arithmetic]
The closure package now contains two complementary certificates: (i) the original floating-point Lean-consumed table with machine-epsilon residual checks, and (ii) an exact symbolic certificate over \(\mathbb Q\). The first is directly wired into the current Lean closure modules; the second is an exact proof-object for strict audit trails.
\end{remark}
For the ``outside-seed-direction'' mechanism, one may choose a witness $X_{k_\star}$ outside the raw seed span (guaranteed by the linear obstruction theorem) and then produce a short Lie word $W_\star$ in $G_2\cup\{\Delta\}$ such that
\[
W_\star=\sum_{k=1}^{28} a_k X_k,\qquad a_{k_\star}\neq 0,
\]
typically appearing after two or three nested brackets in the closure search.

To bridge linear span $\le 15$ to full Lie closure $=28$, the proof uses:
\begin{enumerate}
  \item linear obstruction:
  \path{finrank_span_G2_union_Delta_le_15};
  \item basis directions outside raw seed span:
  \path{exists_so8Generator_not_mem_span_G2_union_Delta};
  \item exhaustive commutator closure and dimension witness:
  \path{lieBracket_in_span} and \path{G2_plus_Delta_closes_to_so8}.
\end{enumerate}
So the additional directions are not assumed: they are produced by iterated bracketing and certified by the full coefficient table.

\section{Causal sets and the Spin(8) closure pipeline}
\label{sec:cst-forcing}

This section records how causal set theory \emph{motivates} the same growth-to-phase-to-algebra narrative, while keeping the logical dependencies explicit. A locally finite partial order does not, by itself, determine an octonion product, a color axis, or a phase plane in \(\mathbb{R}^8\); those inputs match the HQIV/HQVM package imported from the companion development.

\begin{theorem}[Causal-set compatibility with the \(\mathfrak{so}(8)\) closure pipeline]
\label{thm:cst-so8}
Suppose a discrete shell-counting function \(A_C(m)\) satisfies the same increment law as Proposition~\ref{prop:growth-identities},
\[
A_C(m+1)-A_C(m)=8(m+2),
\]
and fix the curvature density exponent \(\alpha>0\) (for concreteness \(\alpha=3/5\) as in HQIV). Then the constructions of Sections~\ref{sec:growth}--\ref{sec:rapidity} define a cumulative channel \(K_C\), a normalized readout \(\Omega_C\), and a phase map--driven observable \(\mathcal{R}_C\) exactly as for the abstract index \(m\).

Assume in addition the \emph{same external Lie data} as in Section~\ref{sec:closure}: the standard matrix realization of the octonionic derivation algebra \(G_2\subset\mathfrak{so}(8)\) and the HQVM/Fano-aligned phase-lift generator \(\Delta_C=\Delta\) on the fixed \((e_1,e_7)\) plane. Then
\[
\langle G_2,\Delta_C\rangle_{\mathrm{Lie}}=\mathfrak{so}(8).
\]
In other words, a causal set (or any other discrete model) whose shell increments match the HQIV quadratic law \emph{equips that modeling choice} with the same analytic curvature channel and, once the octonionic seed is imported, the same certified Lie closure. No uniqueness claim is made that CST axioms single out this \(\Delta\) or \(G_2\) among all Lie data compatible with the growth rate.
\end{theorem}

\begin{proof}
The analytic conclusions are identical to Theorem~\ref{thm:main-curvature} and Section~\ref{sec:rapidity} because they depend only on the discrete increment and \(\alpha>0\). The Lie closure is the same certificate as Theorem~\ref{thm:full-lie-closure}: its proof uses only bracket relations among the fixed matrices \(G_2\cup\{\Delta\}\), not the ontological interpretation of the index \(m\) as a causal-set cardinality.
\end{proof}

\paragraph{Same interpretive caveat.}
The phase readout \(\mathcal{R}_C\) inherits the logic of Section~\ref{sec:rapidity}: it is \emph{not} claimed to be classical rapidity in a continuum Lorentzian limit, nor a Berry phase without further structure.

\begin{remark}[Higher-order gauge group]
In the continuum limit the emergent \(\mathrm{Spin}(8)\) organizes triality automorphisms of the Clifford algebra \(\mathrm{Cl}(8)\) and supplies the gauge structure for internal degrees of freedom (color, generation, etc.) while the spacetime geometry remains governed by the underlying causal partial order. This separation of ``spacetime'' and ``internal'' symmetries is precisely what the growth-to-closure mechanism enforces: the polynomial degree of the shell growth (here quadratic, corresponding to \(3{+}1\) dimensions in the HQIV null-shell reading) selects the octonionic seed used in this paper, which in turn yields the \(28\)-dimensional closure once that seed is adjoined to the phase lift. The same \emph{template} extends to other \(n\) via Fourier twiddles on \(S^{n-1}\); carrying the exceptional \(G_2\) sector requires the octonionic \(8\)-dimensional carrier as in Section~\ref{sec:closure}.
\end{remark}

\begin{remark}[Why the octonionic seed is intrinsic to \(3{+}1\)D causal growth (narrative synthesis)]
The critique that the octonionic seed \(G_2 \cup \{\Delta\}\) must be ``supplied'' is formally correct at the level of the Lie-algebra statement. However, within the present framework the seed is not arbitrary: it is selected by the \(3{+}1\) dimensionality of the causal growth model used here. The quadratic shell-counting law \(A(m) \sim m^2\) encodes three spatial dimensions (null shells transverse to a light-cone in \(3{+}1\) spacetime in the HQIV reading). In three dimensions the natural rotation algebra \(\mathfrak{so}(3)\) admits a spin lift \(\mathrm{Spin}(3) \simeq \mathrm{SU}(2)\), but to obtain an exceptional structure whose triality automorphism group \(G_2\) stabilizes a single ``color'' axis while acting irreducibly on the remaining seven imaginary units, one passes to the octonion algebra \(\mathbb{O}\) whose imaginary part is seven-dimensional. The Fano-plane multiplication table of the octonions is the combinatorial object used in HQVM to coordinate those seven directions with the \(3\)D null lattice bookkeeping. The phase-lift generator \(\Delta\) (mixing the electromagnetic direction \(e_1\) with the \(G_2\)-stabilized color axis \(e_7\)) is the HQIV/HQVM choice compatible with both the null-shell growth package and the requirement that the internal algebra close under brackets to \(\mathfrak{so}(8)\). In lower-dimensional analogs the same mechanism yields only the quaternionic closure \(\mathfrak{so}(4)\) (the \(\mathfrak{so}(3)+\Delta_4\) toy model in Section~\ref{sec:closure}); the jump to the \(28\)-dimensional exceptional algebra is the octonionic certificate proved in this repository. Separately from the discrete shell index \(m\in\mathbb{N}\), the classical Hurwitz fact---finite-dimensional normed division \(\mathbb{R}\)-algebras occur only in dimensions \(1,2,4,8\)---motivates why an \(8\)-dimensional real carrier is a natural place to house internal rotations once one commits to a division-algebra narrative; it is \emph{not} a corollary of polynomial growth on \(\mathbb{N}\), nor a literal ``continuation in \(m\)'' through \(\mathbb{R}\to\mathbb{C}\to\mathbb{H}\to\mathbb{O}\). Turning such algebra classification into a uniqueness theorem from CST axioms (or from shell counting alone) is left to future work.
\end{remark}

\begin{example}[Minimal discrete illustration: first two shells and an outside-span bracket]
Take the normalization index \(m_\ast=1\), so \(K(1)=\rho(1)=1\) and \(\Omega(1)=1\). At shell index \(n=2\),
\[
K(2)=\rho(1)+\rho(2)=1+\frac{1+\alpha\log 2}{2},
\]
which for \(\alpha=\tfrac35\) is \(1+\tfrac12\bigl(1+\tfrac35\log 2\bigr)\approx 1.708\). Thus \(\Omega(2)=K(2)/K(1)\approx 1.708\). With the concrete phase map \(\mathcal F(\phi,t,\Omega)=\Omega\cdot\mathcal H(\phi,t)\) and unit base phase \(\mathcal H(\phi,t)=1\) for illustration, the normalized readout jumps from \(\mathcal R(\phi,t,1)=1\) to \(\mathcal R(\phi,t,2)\approx 1.708\), so the phase increment \(\theta(2)=\mathcal R(\phi,t,2)-\mathcal R(\phi,t,1)\approx 0.708\) multiplies the fixed generator \(\Delta\) in
\[
U(\varepsilon,2)=\exp(\varepsilon\,\theta(2)\,\Delta),\qquad \Delta_{1,7}=1,\ \Delta_{7,1}=-1.
\]
Independently of this numeric illustration, the Lie closure certificate already contains a two-bracket word whose expansion uses a direction outside the raw \(15\)-dimensional linear span of \(G_2\cup\{\Delta\}\): for the numbered basis \(\{X_k\}_{k=0}^{27}\) in Lean,
\[
[X_0,X_1]=\sum_{k=0}^{27} c_{01}^k X_k,
\qquad
c_{01}^{2}=\frac{-13165551793093}{100000000000000}\neq 0,
\]
exhibiting a nonzero coupling to a basis direction \(X_2\notin \mathrm{span}(G_2\cup\{\Delta\})\). (The companion coefficient \(c_{01}^{27}\) recorded in Section~\ref{sec:closure} is likewise nonzero.) This is the same phenomenon as \([J_{12},J_{14}]=J_{24}\) in the \(\mathfrak{so}(4)\) toy model: closure directions are generated by brackets, not assumed in the linear seed.

\smallskip
\noindent\textit{Narrative reading.} In the HQIV synthesis, already at low shell index the quadratic growth law, the \(8\)-dimensional mode unlock, and the \(G_2\) companions work together so that the phase-lift \(\Delta\) participates in Lie words that open directions outside the raw \(15\)-dimensional span; the octonionic table is the concrete certificate carrier in which that opening is realized machine-checkably.
\end{example}

\begin{remark}[Machine-checked causal-set package]
The Lean theorem \texttt{Hqiv.Story.causal\_set\_growth\_forces\_spin8} is definitionally the same packaged statement as \texttt{causal\_rapidity\_forces\_octonion}: it composes the shared-manifold common phase-readout transport lemma, the curvature lower bound and divergence theorems for \texttt{curvature\_integral}, the reference-shell normalization \texttt{omega\_k\_partial\_at\_reference}, and the \texttt{so8Generator} span witness \emph{supplied as the same explicit hypothesis} as in the causal-rapidity package (the closure certificate itself remains the machine-checked \path{G2_plus_Delta_closes_to_so8} table). The discrete shell increment identity matching \(A(m+1)-A(m)=8(m+2)\) is \texttt{Hqiv.Story.causal\_shell\_growth\_law} (alias of \texttt{new\_modes\_succ} in \path{Hqiv/Geometry/OctonionicLightCone.lean}); the transport step is \texttt{Hqiv.Story.rapidity\_from\_causal\_poset} (alias of \texttt{shared\_manifold\_induces\_common\_rapidity}).
\end{remark}

\section{Methodological Unification}

\begin{proposition}[Shared rigid-output pattern]
\label{thm:forced-structure}
Assume:
\begin{enumerate}
  \item a uniform polynomial growth law on a discrete index space, with \(\alpha>0\), producing the channel \(K\) as in Theorem~\ref{thm:main-curvature}; and
  \item an octonionic seed \(G_2\cup\{\Delta\}\subset \mathfrak{so}(8)\) with the certified Lie-closure data from Section~\ref{sec:closure}.
\end{enumerate}
Then three structural outputs follow from the listed hypotheses:
\begin{itemize}
  \item a sign-flip-invariant normalized cumulative phase readout \(\mathcal R\) (``rapidity'' only by interpretive convention);
  \item exact \(2\pi\)-periodic holonomy on the distinguished phase plane;
  \item upon adjoining a suitable seed algebra, full Lie closure to \(\mathfrak{so}(n)\) (specialized in this paper to \(\mathfrak{so}(8)\)).
\end{itemize}
\end{proposition}

\begin{proof}
The first conclusion is exactly the content of Theorem~\ref{thm:main-curvature}, the normalization lemma for \(\Omega\), and Definition~3.3 of the phase readout \(\mathcal R\). The second conclusion is Theorem~\ref{thm:universal-holonomy}. The third conclusion is Theorem~\ref{thm:full-lie-closure} together with the linear-vs-Lie distinction made explicit by the obstruction lemma.
All three are ``rigid-output'' statements under the listed hypotheses: divergence and positivity yield a nontrivial phase readout \(\mathcal R\); the phase-lift data yield exact planar \(2\pi\)-periodicity (Theorem~\ref{thm:universal-holonomy}); and iterated commutators complete the raw seed to the full \(28\)-dimensional algebra in the octonionic certificate (Theorem~\ref{thm:full-lie-closure}).
\end{proof}

\begin{remark}[Template, not causal physical derivation]
Proposition~\ref{thm:forced-structure} is methodological: it does not assert that physical growth data \emph{derives} \(\mathfrak{so}(8)\) (or conversely). The third bullet---full Lie closure in the octonionic seed case---uses only the certified \(G_2\cup\{\Delta\}\) data and bracket algebra; it does not depend on the numerical value of \(\alpha\), on the specific polynomial \(A(m)\), or on HQIV axioms, once \(\Delta\) is the same matrix used in Section~\ref{sec:closure}. The first two bullets use positivity and divergence of \(K\) and the phase-map axioms for \(\mathcal R\). Thus the \emph{closure certificate} and the \emph{growth-to-rapidity chain} can be read as parallel instantiations of a common proof pattern, not as a single causal physical implication from mode counting to exceptional Lie symmetry.
\end{remark}

\subsection*{Causal rapidity package (machine-checkable form)}
For the null-shell growth law
\[
A(m)=4(m+2)(m+1),
\]
combine:
\begin{enumerate}
  \item the optional shared-manifold transport interface (Appendix~\ref{app:transport}; single ambient phase channel seen from both embeddings);
  \item the curvature/divergence chain from Section~2 (\(K(n)\ge H_n\), \(K\to+\infty\), \(\Omega\)-normalization at reference shell);
  \item the certified \(G_2+\Delta\) closure package from Section~5.
\end{enumerate}
Then the full conjunction is packaged in Lean as one theorem:
\[
\texttt{Hqiv.Story.causal\_rapidity\_forces\_octonion}.
\]
Concretely, this theorem returns:
\begin{itemize}
  \item existence of a common rapidity observable on the shared manifold interface;
  \item positivity at reference shell and divergence of the curvature integral;
  \item reference normalization \(\Omega(\text{reference})=1\);
  \item the imported \(28\)-dimensional \(\mathfrak{so}(8)\) closure witness.
\end{itemize}
So the packaged conjunction is machine-checkable in Lean under the hypotheses stated above.

\begin{theorem}[Stacked spin-channel package in \texorpdfstring{$1,2,3$}{1,2,3} dimensions (Lean form)]
Assume the \(d=3\) shell-count growth law
\[
A(m)=4(m+2)(m+1),
\]
a shared-manifold causal transport bridge with common phase readout (same interface as Appendix~\ref{app:transport}), and the certified \(\mathfrak{so}(8)\) closure witness from Section~5 as interface hypothesis. Then:
\begin{enumerate}
  \item a common rapidity observable exists on the shared manifold interface (spin-channel compatibility at each layer);
  \item the curvature channel at the reference shell is positive and the cumulative channel diverges, so the normalized phase readout \(\Omega(n)\) is unbounded as \(n\to\infty\);
  \item the reference normalization satisfies \(\Omega(\mathrm{reference})=1\), and the phase-lift action provides exact \(2\pi\)-periodic holonomy on the distinguished plane;
  \item the imported \(G_2+\Delta\) closure package provides the \(28\)-dimensional \(\mathfrak{so}(8)\) (octonionic-layer) witness.
\end{enumerate}
Interpretation: the \(1\)D/\(2\)D/\(3\)D spin-channel constraints stack, planar holonomy closes at \(2\pi\), and the same closure certificate supplies octonionic-layer structure once the \(G_2\cup\{\Delta\}\) hypothesis is imported; this is a packaging statement, not a uniqueness theorem for discrete spacetimes.
\end{theorem}

\begin{proof}
This is exactly the theorem \texttt{Hqiv.Story.causal\_rapidity\_forces\_octonion} (equivalently \texttt{causal\_growth\_forces\_real\_imaginary\_layers\_d3}), which composes the already proved components
\texttt{shared\_manifold\_induces\_common\_rapidity},
\texttt{curvature\_integral\_ref\_pos},
\texttt{curvature\_integral\_tends\_to\_atTop},
\texttt{omega\_k\_partial\_at\_reference},
and the imported closure witness hypothesis.
\end{proof}

\begin{remark}[3D associativity-breaking interface and observable superposition]
At the manuscript level, the \(d=3\) octonionic layer is interpreted as allowing mild associativity breaking in composition order. In the current Lean packaging, this is represented conservatively as an interface implication
\[
\text{(mild 3D associativity breaking)}\Longrightarrow \text{(superposition-form observable)},
\]
without introducing new octonion-multiplication axioms into the verified core. This split is formalized as
\[
\texttt{Hqiv.Story.stacked\_spin\_channel\_forcing\_d3\_with\_superposition\_interface}.
\]
So the proved geometric/algebraic package is unchanged, while the 3D superposition reading is attached as an explicit hypothesis-level bridge.
\end{remark}

\begin{remark}[Corollary interface: emergent handedness]
Within this package, once an orientation rule for \(\Delta\) is fixed and reflection pairing \([x,|x|]\) is imposed, chirality selection statements (CISS-style handedness bias) reduce to an additional transport-to-material map. In other words, the geometric/algebraic packaging theorem is complete; the CISS claim is an interface corollary requiring one extra domain-specific bridge.
\end{remark}

\section{Foundational and Adjacent Literature}

This paper touches five classical lines of work.

\subsection*{Octonions, \texorpdfstring{$G_2$}{G2}, triality, and \texorpdfstring{$\mathfrak{so}(8)$}{so(8)}}
The octonionic derivation algebra $G_2$ and the triality structure of $\mathrm{Spin}(8)$ are standard foundations for exceptional and rank-8 geometry. Classical references include Cartan's work on exceptional Lie groups and modern expositions such as Baez's survey on octonions \cite{baez2002} and Springer--Veldkamp \cite{springer2000}. The present contribution is not a classification result, but a closure pipeline with an explicit \emph{linear-vs-Lie} distinction and machine-checked closure certificate.
On the physics side, \(\mathrm{SO}(8)\) spinors and triality appear in the construction of space-time supersymmetry and GSO projections in ten-dimensional superstring theory (see \cite{green1987,polchinski1998}); compactification and anomaly constraints further tie low-energy gauge groups to the same representation-theoretic backbone. More recent work embeds gauge structure and three generations via division algebras and related octonionic constructions \cite{furey2018,borsten2014}, providing broader motivation for treating \(G_2\)-to-\(\mathfrak{so}(8)\) closure as physically meaningful structure beyond this paper's formal scope.
This derivation also aligns with discretization programs emphasizing emergence of exceptional structure from combinatorial data (spin-foam/causal-set viewpoints), where rapidity acts as a first-order bridge from growth channels to orthogonal-group symmetry organization \cite{baez2002}.

\subsection*{Growth-implies-structure paradigm}
The statement that growth laws force geometric or algebraic structure is reminiscent of Gromov's polynomial-growth rigidity paradigm \cite{gromov1981}. In discrete and combinatorial geometry, curvature notions of Forman and Ollivier \cite{forman2003,ollivier2009} similarly connect local transport and growth behavior to global geometric constraints. Our curvature channel is a direct constructive instance of this philosophy: the chosen growth law yields an explicit positive, divergent cumulative curvature functional.

\subsection*{Causal set gauge fields and sequential growth}
Prior CST literature has pursued two main routes to gauge structure. Sverdlov (arXiv:0807.2066, 2008) and follow-up works express continuum Yang--Mills Lagrangians in terms of CST holonomies, causal relations, and volumes, effectively embedding gauge fields onto a pre-existing causal set. Classical sequential growth (CSG) models (Rideout and Sorkin, Phys.\ Rev.\ D \textbf{61}, 024002, 2000) supply the dynamical framework in which causal sets accrete element by element; labels on the growth order are pure gauge and covariance is imposed by hand. These approaches are largely post-hoc: the gauge group is introduced after the causal set exists, or recovered only in a continuum limit. Foundational references for the causal-set program include \cite{bombelli1987,sorkin2005}. The present manuscript instead adopts the HQIV null-shell counting law as a fixed analytic input and packages the curvature/rapidity/closure chain of Sections~\ref{sec:growth}--\ref{sec:closure} accordingly; Section~\ref{sec:cst-forcing} states the causal-set interpretation at the level of hypotheses used in Lean.

\subsection*{Rapidity and geometric phase viewpoints}
Rapidity is classically a hyperbolic parameter in Lorentzian kinematics; geometric phase frameworks generalize phase observables to broader manifold settings. The novelty here is the derivation chain
\[
\text{growth}\to K\to \Omega\to \mathcal R,
\]
which promotes rapidity from a coordinate choice to an induced index-space observable attached to a normalized curvature channel.

\subsection*{Formal verification context}
Lean/mathlib is increasingly used for algebra and geometry formalization. This manuscript adopts a derivation-first style in the main text and records the nontrivial formal artifacts in Appendix~A, so claims can be audited at proof-object level without turning the paper into a code index.
For submission hygiene, all cited software artifacts (repository URLs, scripts, and generated certificates) should be evaluated at stable public snapshots/DOIs so referees can reproduce the exact reviewed state; the archive \PaperArtifactBundleZip{} frozen for \path{papers/closure.tex} is the frozen bundle intended for that purpose (this long-form file is not part of that archive).

\section{Conclusion}

The derivation chain established in this paper is:
\[
\text{uniform growth}\Longrightarrow K\Longrightarrow \Omega\Longrightarrow \mathcal R
\Longrightarrow \Delta
\Longrightarrow \mathfrak{so}(8).
\]
When a discrete model's shell increments match the HQIV quadratic law, Theorem~\ref{thm:cst-so8} records how the same analytic channel and the \emph{imported} \(G_2\cup\{\Delta\}\) Lie data fit together; it does not assert that causal-set axioms uniquely determine that Lie data.
Historically, the long geometric route and the concise \(G_2+\Delta\) route are the same construction seen at two resolutions: the latter is the algebraic distillation of the former. What is proved here is the full derivation flow in a compact form, with closure certificates and transport components audited in Lean.
Interpretively, the derivation chain is used here as a physics-motivated structural map (light-cone/discrete-shell setting and causal-set motivation in Section~\ref{sec:cst-forcing}), not as a claim that growth data alone physically causes \(\mathfrak{so}(8)\) without the explicit Lie data of Section~\ref{sec:closure}. The same rigid-output template can be transplanted to other domains where cumulative growth observables are central.
Although the primary result concerns the octonionic case, the same mechanism appears in a direct 6-dimensional analog in \(\mathfrak{so}(4)\): scalar channel \(\rightarrow\) phase-lift seed \(\rightarrow\) Lie closure. The identical symbolic certificate pipeline provides an exact rational structure-constant table for the quaternionic seed \((\mathfrak{so}(3)+\Delta_4)\), and the corresponding artifacts are included in \PaperArtifactBundleZip{} frozen for \path{papers/closure.tex}.

\section{Outlook (programmatic suggestions only)}
\label{sec:outlook}
\paragraph{Scope (not part of the proved core).}
Everything below is \emph{programmatic}: possible directions for future work, stated without quantitative claims or new axioms. No prediction (e.g.\ a fixed twiddle multiplier at low \(n\), or a closed value for \(\zeta(3)\)) is asserted as a theorem here; any such statement would belong to a separate paper with its own hypotheses and verification.
The bullet list is intentionally speculative and is \emph{not} used anywhere in Sections~\ref{sec:growth}--\ref{sec:closure}:
\begin{itemize}
  \item extension of the reflection-pairing mechanism to full critical-line constraints for zeta-type zero sets;
  \item systematic acceleration schemes for constants (e.g., \(\pi\)) from the same twiddle/closure certificate stack in higher dimension;
  \item broader geometric interpretations (including holographic-style dual descriptions) driven by the same growth-to-holonomy template.
\end{itemize}
These are deferred to follow-up work; the present paper establishes only the proved core chain and its machine-checked closure/certificate components.

\subsection*{Schematic chain}
\[
\text{growth}\Longrightarrow K\Longrightarrow \Omega\Longrightarrow \mathcal R
\Longrightarrow \Delta\Longrightarrow \text{Lie words}\Longrightarrow \mathfrak{so}(n)\text{ closure},
\]
generating downstream exact certificates and hypercomplex Fourier twiddles on \(S^{n-1}\).

The same normalized phase readout could be compared, informally, to classical rapidity on hyperbolic surfaces, to geometric/Berry phases, and to critical-line symmetry heuristics for zeta-type functions via index/reflection pairings; none of these bridges are proved in the present manuscript. Hypercomplex Fourier analysis on \(S^{n-1}\) and accelerated-constant programs are likewise mentioned only as possible extensions.

\begin{remark}[Holonomy as parallel transport]
In a smooth gauge theory, parallel transport along a loop yields a holonomy \(U=\mathcal P\exp\oint A\); curvature is the obstruction to trivial holonomy for infinitesimal loops (Ambrose--Singer). Here, in the discrete layer, the phase-lift generator \(\Delta_{ij}\) exponentiates to \(U(\varepsilon,n)=\exp(\varepsilon\,\theta(n)\,\Delta_{ij})\), and Theorem~\ref{thm:universal-holonomy} gives \(\exp(2\pi\,\Delta_{ij})=I\): a \emph{closed} increment in the rapidity parameter closes the loop holonomy on the embedded plane. The analogy is direct: nontrivial \(\theta(n)\) produces rotation in the plane; requiring \(2\pi\)-periodicity is the discrete shadow of demanding consistent transport around a closed phase cycle. For evolving systems where cumulative mode counting drives a divergent \(K(n)\), the same primitive phase channel can be read as encoding an effective transport memory; any ``bulk'' or holographic language in Section~\ref{sec:outlook} remains interpretive and is not used in the formal proofs.
\end{remark}

\appendix
\section{Formal Verification Map}

The text above is derivation-first. Corresponding machine-checked artifacts are:

\paragraph{Minimal Lean build (repository).}
Unpack \PaperArtifactBundleZip{} (submission bundle for \path{papers/closure.tex}) (or clone the live \cite{hqiv-lean} repository at a matching commit). From the directory containing \path{lakefile.toml}, run \texttt{lake build HQIVPaperClaims}. In the current tree this target is intentionally lightweight: it checks the growth/curvature chain, SAT rapidity transport, and causal forcing package, and imports the symbolic SO(8) closure interface (\path{Hqiv/SO8ClosureSymbolic.lean}) rather than elaborating the heavy matrix-certificate closure cone.
For the full certified closure build (Lie-bracket coefficient tables and closure proofs), run the dedicated target \texttt{lake build HQIVSO8Closure}.
Separately, the optional strong-color target \texttt{lake build HQIVStrongColorSu3Certificate} elaborates the generated \(\mathfrak{su}(3)\) \(f^{abc}\) \texttt{@[simp]} table used by the triplet chart modules (\path{Hqiv/Physics/StrongColorSu3fStructureSimp.lean}); it is intentionally outside the default \texttt{HQIVLEAN} import cone.

\subsection*{Growth and curvature}
\begin{itemize}
\item \texttt{alpha\_eq\_3\_5}, \texttt{latticeAlphaRatio\_eq\_alpha}
\item \texttt{curvature\_integral\_ge\_harmonic}
\item \texttt{curvature\_integral\_tends\_to\_atTop}
\item \texttt{omega\_k\_partial\_at\_reference}
\end{itemize}

\subsection*{Rapidity channel and transport interface}
\begin{itemize}
\item \texttt{rapidity\_tip\_recovers\_harmonic\_step}
\item \texttt{shared\_manifold\_induces\_common\_rapidity}
\item \texttt{shares\_threshold\_of\_pointwise\_bounds}
\item \texttt{shares\_threshold\_of\_step\_bounds}
\item \texttt{causal\_growth\_forces\_real\_imaginary\_layers\_d3}
\item \texttt{causal\_rapidity\_forces\_octonion}
\item \texttt{causal\_set\_growth\_forces\_spin8} (definitional alias of \texttt{causal\_rapidity\_forces\_octonion})
\item \texttt{causal\_shell\_growth\_law} (discrete increment \(A(m+1)-A(m)=8(m+2)\) for \texttt{available\_modes})
\item \texttt{rapidity\_from\_causal\_poset} (alias of \texttt{shared\_manifold\_induces\_common\_rapidity})
\item \texttt{stacked\_spin\_channel\_forcing\_d3\_with\_superposition\_interface}
\end{itemize}

\subsection*{\texorpdfstring{$\mathfrak{so}(8)$}{so(8)} closure}
\begin{itemize}
\item lightweight symbolic interface module:
\texttt{Hqiv/SO8ClosureSymbolic.lean}
\item \texttt{so8\_generators\_linear\_independent}
\item \texttt{lieBracket\_in\_span}
\item \texttt{span\_so8Generators\_eq\_skewAdjointOne}
\item \texttt{finrank\_span\_G2\_union\_Delta\_le\_15}
\item \texttt{exists\_so8Generator\_not\_mem\_span\_G2\_union\_Delta}
\item \texttt{G2\_plus\_Delta\_closes\_to\_so8}
\item symbolic certificate generator (in \PaperArtifactBundleDir/):
\texttt{scripts/generate\_symbolic\_lie\_closure\_certificate.py}
\item exact proof-object artifact:
\texttt{artifacts/so8\_symbolic\_certificate.json}
\item implementation note: the generator \texttt{scripts/generate\_symbolic\_lie\_closure\_certificate.py} is \emph{specialized} to the HQVM octonionic \(8\times8\) seed (rank-\(28\) \(\mathfrak{so}(8)\) closure). The Fourier/holonomy layer in the main text is stated for general embedded planes \(\Delta_{ij}\in\mathfrak{so}(n)\); a separate script would be needed to emit full rational certificates for other \(n\).
\item toy-model symbolic generator:
\texttt{scripts/generate\_symbolic\_so4\_certificate.py}
\item toy-model exact artifact:
\texttt{artifacts/so4\_symbolic\_certificate.json}
\end{itemize}

\paragraph{One-line external audit.}
In Python: \texttt{import json}, then \texttt{d = json.load(open("artifacts/so8\_symbolic\_certificate.json"))} (path relative to \PaperArtifactBundleDir/ after unpacking \PaperArtifactBundleZip{} for \path{papers/closure.tex}); inspect keys for basis labels and rational structure constants. In computer algebra, map numerator/denominator pairs to \(\mathbb Q\) and spot-check a bracket against \([X_i,X_j]=\sum_k c^k_{ij}X_k\). Lean consumes the floating-point table in \path{Hqiv/GeneratorsLieClosureData*.lean} when building \texttt{HQIVSO8Closure}; the JSON is an independent exact-arithmetic certificate for the same algebra.

\section{Supplements for readers and formal auditors}

\subsection{Optional abstract transport interface}
\label{app:transport}
The following diagram is the same optional interface formerly placed in Section~\ref{sec:rapidity}; it is not used in the discrete core proofs. Let $X$ be a carrier set (or space) with maps $e_1:X_1\to X$ and $e_2:X_2\to X$ and a shared observable $r:X\to\mathbb R$. Then
\[
r_1:=r\circ e_1,\qquad r_2:=r\circ e_2
\]
are induced channels on each domain. Any bound on $r$ over one image transports to the other channel under the same threshold hypotheses.
\[
\begin{CD}
X_1 @>{e_1}>> X @<{e_2}<< X_2\\
@V{r_1}VV @V{r}VV @VV{r_2}V\\
\mathbb R @= \mathbb R @= \mathbb R
\end{CD}
\]

\subsection{HQIV in one slide}
\label{app:hqiv-box}
\emph{Horizon-Quantized Informational Vacuum} (HQIV) is a companion physics framework in which overlapping causal horizons and entanglement monogamy motivate discrete null-shell mode counting. This paper \emph{imports} three concrete pins from that story: (i) the quadratic shell law \(A(m)=4(m+2)(m+1)\) and increment \(8(m+2)\); (ii) the curvature imprint exponent \(\alpha=3/5\) tied to Brodie's horizon thermodynamics \cite{brodie2026} and the lattice-ratio identity in Lean; (iii) the HQVM octonion/Fano realization of \(G_2\) and the phase-lift \(\Delta\) in the \((e_1,e_7)\) plane. None of these pins are reconstructed here from causal-set axioms alone; they are the explicit hypotheses that make the Lean modules \path{Hqiv/Geometry/OctonionicLightCone.lean}, \path{Hqiv/Geometry/SATRapidityManifold.lean}, and the Lie-closure tables under \path{Hqiv/GeneratorsLieClosureData*.lean} mutually interoperable (all included in \PaperArtifactBundleZip{} frozen for \path{papers/closure.tex}).

\subsection{Lean dependency sketch (ASCII)}
\label{app:lean-deps}
\begin{verbatim}
OctonionicLightCone  (growth, rho, K, Omega, curvature_integral, referenceM)
        |
        +--> SATRapidityManifold (shared_manifold_induces_common_rapidity, ...)
        |
        v
Hqiv.Story.CausalRapidityForcing
   (causal_rapidity_forces_octonion  [= causal_set_growth_forces_spin8])
        |
        +--> [hypothesis: so8Generator span witness dim 28]
        |
        v
HQIVPaperClaims --> Hqiv/SO8ClosureSymbolic  (lightweight symbolic closure interface)
HQIVSO8Closure / GeneratorsLieClosure        (full certified G2_plus_Delta_closes_to_so8)
\end{verbatim}

\section*{Acknowledgments}
AI-assisted drafting, editing, and formal-verification tooling were used during manuscript preparation (including Composer, Grok, and related coding assistants). Final mathematical claims, Lean artifacts, and script outputs were reviewed and validated by the author.

\begin{thebibliography}{99}
\bibitem{hqiv-lean}
HQIV-LEAN repository, \url{https://github.com/disregardfiat/hqiv-lean}, 2026; reproducibility subtree \PaperReproBundleDir/ (see \PaperReproBundleManifest) is frozen for \path{papers/closure.tex} inside the release distribution \ReleaseDistArchive{}.
\bibitem{brodie2026}
K. Brodie, \emph{Derivation of Quantised Inertia from Jacobson Thermodynamics}, Zenodo (2026), doi:\,\url{https://doi.org/10.5281/zenodo.18706746}.
\bibitem{baez2002}
J. C. Baez, \emph{The Octonions}, Bull. Amer. Math. Soc. 39 (2002), no. 2, 145--205.
\bibitem{forman2003}
R. Forman, \emph{Bochner's method for cell complexes and combinatorial Ricci curvature}, Discrete Comput. Geom. 29 (2003), 323--374.
\bibitem{gromov1981}
M. Gromov, \emph{Groups of polynomial growth and expanding maps}, Publ. Math. IHES 53 (1981), 53--78.
\bibitem{green1987}
M. B. Green, J. H. Schwarz, and E. Witten, \emph{Superstring Theory, Vols. 1--2}, Cambridge University Press, 1987.
\bibitem{furey2018}
C. Furey, \emph{Three generations, two unbroken gauge symmetries, and one eight-dimensional algebra}, Phys. Lett. B 785 (2018), 84--89.
\bibitem{polchinski1998}
J. Polchinski, \emph{String Theory. Vol. I: An Introduction to the Bosonic String}, Cambridge University Press, 1998.
\bibitem{borsten2014}
L. Borsten, M. J. Duff, L. J. Hughes, S. Nagy, \emph{Magic square from Yang--Mills squared}, Phys. Rev. Lett. 112 (2014), 131601, arXiv:1301.4176 (division algebras and dualities linked to triality-type symmetry data).
\bibitem{ollivier2009}
Y. Ollivier, \emph{Ricci curvature of Markov chains on metric spaces}, J. Funct. Anal. 256 (2009), 810--864.
\bibitem{springer2000}
T. A. Springer and F. D. Veldkamp, \emph{Octonions, Jordan Algebras and Exceptional Groups}, Springer, 2000.
\bibitem{sorkin2005}
R. D. Sorkin, \emph{Causal sets: Discrete gravity}, in \emph{Lectures on Quantum Gravity}, eds.\ A. Gomberoff and D. Marolf, Springer, 2005, pp.\ 305--327.
\bibitem{bombelli1987}
L. Bombelli, J. Lee, D. Meyer, and R. D. Sorkin, \emph{Space-time as a causal set}, Phys. Rev. Lett. 59 (1987), 521--524.
\end{thebibliography}

\end{document}